% SOURCE docs/release/2026-10-29/source/docs/writing-repair/2026-10-29/src/dd7bd35bd9f6df8e.tex CHAR 0
\chapter{Identification, counterfactuals and effect sets}
\label{exec17:s3}

A bank can know that an offer increases average completion and still be unable
to say which completed applications the offer caused. It can identify the
effect of sending an encouragement without identifying the effect of making
every customer receive the service. It can also possess a valid identifying
formula whose numerical evaluation is too fragile to support the proposed
decision. These are different limits on what a causal product can say. The
choice between software components depends on preserving those limits through
the calculation.

This chapter develops that choice around a binary servicing action and a
binary successful outcome. The examples are project-authored probability
laws, evaluated exactly, with a stipulated value of USD 100 per success and
USD 12 per action received. They are neither estimates of bank behavior nor
measurements of upstream software. Their role is to make a proposal testable:
an implementation can be asked to preserve the target, recover the
appropriate answer or set, and explain a refusal without confusing a
scientific obstacle with a computational one.~\citep{exec17_s3_exact}

The argument first separates adjustment, general identification, individual
counterfactuals and encouragement effects. It then develops the full binary
instrument table to show how changing an assumption changes a particular
identified set. The language comparisons translate those distinctions into
symbolic results, search termination and model-specific calculations. The
closing Ananke and DoWhy GCM studies examine two further boundaries: the
scale of a reported effect and the information needed to retain an
individual's noise. Each comparison therefore asks what information a
returned quantity uses, before asking which engine should compute it.

\section{A causal question is relative to information and a model class}
\label{exec17:s3-target}

Let $A\in\{0,1\}$ denote receipt of a specified intervention, $Y\in\{0,1\}$
the outcome at a fixed horizon, and $C$ information measured before receipt.
The intervention includes its delivery version: sending a message and
completing a counseling session are different actions even if both are
recorded as contact. Write $Y(a)$ for the response under receipt $a$,
$P$ for the observable population law and $\mathcal M$ for the admitted
causal models. A common target is
\[
 \theta(M)=E_M\{Y(1)-Y(0)\},\qquad M\in\mathcal M.
\]
Consistency links the observed outcome to the potential response under
the action actually received. Assumptions about delivery versions,
interference, information availability and the horizon determine whether
that link describes the banking question. A graph supplies restrictions
on mechanisms and omitted causes; it does not certify these operational
premises.~\citep{pearl2009,swig2013}

For a fixed observable law, define the compatible model class and effect set by
\[
 \mathcal M(P)=\{M\in\mathcal M:P_M=P\},\qquad
 \Theta(P)=\{\theta(M):M\in\mathcal M(P)\}.
\]
The target is point identified at $P$ if this set is a singleton. A graphical
identification theorem usually establishes that a functional $T$ works for
every admitted compatible law, so that $\theta(M)=T(P_M)$ throughout its
stated domain. Special numerical coincidences at one law are weaker evidence.
Conversely, finding two admissible models with the same $P$ and different
$\theta$ disproves point identification for that law and class. Estimating
$P$ more accurately does not distinguish such models. A smaller scientifically
defensible class or additional intervention information may
do so.~\citep{id2006,pearl2009}

The distinction matters when evaluating a proposal. A held-out likelihood
score assesses a distributional prediction. It cannot establish uniqueness
of a causal answer across models that share that distribution. An optimizer
that returns one fitted model has selected one element; it has not shown that
all compatible elements agree. A prior can produce a posterior over models,
but its concentration has to be interpreted with the restrictions and prior
that supplied the missing information.

There is another boundary: a formula may involve $P(Y\mid A=a,C=c)$ where
$P(A=a\mid C=c)=0$. The observed law then does not supply that conditional
response. A smooth extrapolating regression can return a value, but the value
uses additional restrictions. The proposal should declare whether it restricts
the target population, changes the action, acquires information or adopts
an explicit extrapolation model. A numerical nonzero fitted propensity
does not create observed support.

\section{Adjustment changes the comparison, not the action definition}
\label{exec17:s3-adjust}

Suppose consistency, conditional exchangeability
$Y(a)\mathrel{\perp}A\mid C$ and positive receipt probabilities hold in the
target population. Iterated expectation then gives
\begin{align}
 E\{Y(a)\}
 &=\sum_c E\{Y(a)\mid C=c\}P(C=c) \notag\\
 &=\sum_c E\{Y\mid A=a,C=c\}P(C=c).
 \label{exec17:s3-adjustment}
\end{align}
The first equality standardizes to the chosen population; the second uses
exchangeability and consistency on the supported cells. If the sponsor wants
the treated population, replacing $P(C=c)$ by $P(C=c\mid A=1)$ changes the
target. It is not an innocuous estimation option.

An exact generated law makes the cost of ignoring that distinction visible.
Let $P(C=0)=P(C=1)=1/2$ and stipulate
\[
\begin{array}{c|cc|cc}
 &P(A=1\mid C)&P(A=0\mid C)
 &P(Y=1\mid A=0,C)&P(Y=1\mid A=1,C)\\ \hline
 C=0&1/5&4/5&1/10&1/5\\
 C=1&4/5&1/5&3/5&7/10
\end{array}
\]
This is a stipulated exchangeable-within-$C$ construction, rather than
evidence that a real baseline variable removes every common cause.
Both strata have a receipt effect of $1/10$. Standardization yields
$E\{Y(0)\}=7/20$ and $E\{Y(1)\}=9/20$.
The observed receipt groups have different compositions:
\[
 P(C=1\mid A=1)=4/5,\qquad P(C=1\mid A=0)=1/5.
\]
Their outcome means are $3/5$ and $1/5$, producing a raw contrast of
$2/5$. All cells are positive, so the problem is not an absent observation
cell. Receipt selection overweights the high-success stratum in the
treated group.~\citep{exec17_s3_exact}

For universal receipt versus no receipt, define incremental expected
net value $V=100\theta-12$ in USD per eligible customer. The causal contrast
implies $V=-2$; substituting the association implies $V=28$.
A product that chooses an accurate predictive model but compares its
predictions over the wrong receipt groups can therefore recommend the
opposite action. The USD 30 discrepancy here is exact under the generated
law, not a forecast of attainable financial improvement.

The relevant diagnostics follow the derivation. Check treatment versions
and baseline timing; inspect supported cells and overlap; state the
standardization weights; and compare the proposed graph against substantive
common causes and selection. Balance after weighting can diagnose a
particular imbalance in measured variables. It does not demonstrate
exchangeability with respect to unmeasured variables.

Adjustment can also answer the wrong causal question when it conditions
on a mediator. Consider the separate linear structural example
\[
 M=A+\varepsilon_M,\qquad
 Y=2A+3M+\varepsilon_Y,
\]
with randomized $A$ and independent zero-mean errors. Substitution gives
$E\{Y\mid do(A=a)\}=5a$ up to the zero baseline intercept. The total
mean effect of changing $A$ from zero to one is $5$. Holding $M$ fixed
by intervention instead gives a controlled direct effect of $2$.
These are continuous illustrative responses, not the binary success law
above. Including the mediator has changed the intervention comparison.
In a bank, completion assistance and the downstream information it
elicits may both contribute to the total effect; controlling away the
information channel is appropriate only if that direct-effect question
is actually intended.~\citep{exec17_s3_exact,pearl2009}

\section{Completeness must retain the procedure's domain}

Perkovi\'c and colleagues characterize adjustment in DAGs, CPDAGs,
MAGs and PAGs using amenability, exclusion of a forbidden set and blocking
of proper definite-status noncausal paths. Their Theorem 5 gives an
if-and-only-if result for adjustment; Theorem 7 supplies the equivalent
proper-back-door separation condition. The MAG and PAG treatment in this
source excludes selection bias. These conditions require the graph's
endpoint semantics, not merely its adjacency
matrix.~\citep{perkovic2018}

Completeness here means that the criterion characterizes the existence
of adjustment sets within the stated classes and definition. It does
not mean that failure to find such a set proves failure of every
identification method. The same source explicitly distinguishes that
claim from identification by other functionals. An adapter should
therefore report the attempted route and its domain with its result.
No valid adjustment set in the admitted graph leaves a different
question open from nonidentification in the admitted model class.
A search that stops after small candidate sets is weaker again,
because it may not have exhausted the adjustment
problem.~\citep{perkovic2018}

For a supplied CPDAG or PAG, a successful adjustment set must be valid
under the relevant equivalence-class interpretation. Selecting one
convenient orientation can create an answer that does not survive the
remaining graphs. It is reasonable to ask a domain expert for the
orientation of a particular edge, but that additional restriction must
remain visible in the scientific model. The resulting effect is
conditional on it.

This motivates three substantive alternatives when adjustment fails.
One can seek a more general identifying functional under the same
information, retain a set of effects compatible with the information,
or change the information through an admissible experiment or
measurement. Choosing a convenient regression is a fourth computational
action, but it only answers the causal question if its added
restrictions are justified. The banking consequence should be assessed
before choosing among these routes: a set lying entirely below the
receipt cost threshold can already rule out an action.

\section{Joint and conditional intervention questions need different routes}
\label{exec17:s3-idc}

For a semi-Markovian causal graph, ID identifies unconditional
intervention distributions or returns a graphical nonidentification
witness in its formal domain. The district decomposition and hedge
criterion are scientific statements about that model class.
They are not properties of every exception a software implementation
can throw. Shpitser and Pearl's conditional algorithm, IDC, first
uses action--observation exchange where the graphical condition permits
it. After the maximal eligible exchange, it invokes the appropriate
joint identification problem and normalizes. Their Theorems 6--9
establish the conditional reduction, soundness, completeness and
do-calculus result for the stated
queries.~\citep{id2006}

The order matters. In general it is valid to normalize an identified
$P(Y,Z\mid do(A=a))$ to obtain $P(Y\mid do(A=a),Z=z)$
where the denominator is positive. It is invalid to infer that
the conditional target is nonidentified merely because this particular
joint target is nonidentified. The conditional may cancel a component
that the joint requires. A component that supports only joint-first
normalization is consequently a conservative partial route to
conditional queries.

Consider the graph with $A\to Z\to Y$ and $A\leftrightarrow Z$,
with no hidden common cause involving $Y$. Let $U$ be a fair bit,
$E$ an independent Bernoulli$(1/4)$ bit, and let the outcome be generated
from an independent random input with
\[
 P(Y=1\mid Z=z)=1/5+(3/5)z.
\]
Two structural models compatible with the graph are
\[
\begin{array}{c|cc}
 & A&Z\\ \hline
 M_0&U&U\mathbin{\mathrm{XOR}}E\\
 M_1&U&A\mathbin{\mathrm{XOR}}E
\end{array}
\]
The graph permits an arrow with zero effect in a particular model;
no faithfulness assumption is imposed on this witness class.
Observationally $A=U$, so both models give the same strictly positive
law for every triple $(A,Z,Y)$. Under $do(A=a)$, however, $U$ remains
random in $M_0$, whereas $Z$ responds to the set action in
$M_1$.~\citep{exec17_s3_exact}

The interventional means are
\[
\begin{array}{c|cc|c}
 &E\{Y\mid do(A=0)\}&E\{Y\mid do(A=1)\}&\text{contrast}\\ \hline
 M_0&1/2&1/2&0\\
 M_1&7/20&13/20&3/10 .
\end{array}
\]
Thus the observational law cannot identify the marginal effect in
this class; it also cannot identify the joint law of $(Y,Z)$ under
the action. Yet throughout the stated graph class the outcome mechanism
has only the observed parent $Z$ and independent exogenous input.
Consequently, on supported conditioning events,
\[
 P(Y=1\mid do(A=a),Z=z)=P(Y=1\mid Z=z)
                     =1/5+(3/5)z.
\]
The equality follows from the graph's mechanism restriction; agreement
of the two numerical witnesses alone would not prove identification.
All relevant conditioning events are positive in the displayed
witnesses.~\citep{id2006,exec17_s3_exact}

In a servicing interpretation, $Z$ could be a completed information
step and $Y$ a later successful resolution. Knowing the outcome response
among a specified post-action $Z$ group does not tell the bank how much
an action changes the frequency of that group. The conditional success
probability therefore does not substitute for the total-value question.
Conditioning on the observed $Z$ also does not identify the effect
within the principal stratum that would attain that $Z$ under either
action. That stratum is a different, cross-world target.

\section{A population response is not an individual counterfactual}
\label{exec17:s3-coupling}

Single-world intervention graphs make counterfactual variables under
one intervention explicit and connect their factorization to the
extended g-formula. Richardson and Robins distinguish the FFRCISTG
model from the stronger NPSEM with independent errors; single-world
and cross-world restrictions should not be silently interchanged.
Pearl's structural definition computes an individual's counterfactual
by changing the intervention while retaining the same exogenous state.
That gives a counterfactual meaning within a fully specified structural
model. It does not show that observational or experimental marginals
identify that structural specification.~\citep{swig2013,pearl2009}

The distinction is exact even for a randomized action. Suppose the
experiment identifies
\[
 P\{Y(0)=1\}=2/5,\qquad P\{Y(1)=1\}=3/5.
\]
Let $t=P\{Y(0)=1,Y(1)=1\}$. Every response-type law below has these
same two marginals:
\[
\begin{array}{c|cccc}
(Y(0),Y(1))&(0,0)&(0,1)&(1,0)&(1,1)\\ \hline
P&t&3/5-t&2/5-t&t .
\end{array}
\]
Nonnegative cell masses require and suffice for $0\le t\le2/5$.
Randomizing $A$ independently of the response type then produces
the same complete observable $(A,Y)$ law for every $t$.
Intermediate $t$ values are feasible as well as the endpoints,
so these are sharp ranges within the stipulated response-type
class.~\citep{exec17_s3_exact}

The average effect is always $1/5$, but the fraction helped is
$3/5-t\in[1/5,3/5]$ and the fraction harmed is
$2/5-t\in[0,2/5]$. At $t=0$, sixty percent benefit and forty
percent are harmed. At $t=2/5$, twenty percent benefit and nobody
is harmed. Both populations have the same average effect and
randomized outcome frequencies. Larger replications of the same
parallel-arm experiment do not select between these couplings.

The retrospective attribution question is different again. Among
treated successes, consistency and randomization imply
\begin{align*}
 P\{Y(0)=0\mid A=1,Y=1\}
 &=\frac{P\{Y(0)=0,Y(1)=1\}}{P\{Y(1)=1\}}\\
 &=\frac{3/5-t}{3/5}\in[1/3,1].
\end{align*}
One can therefore identify a positive average effect and still be
unable to assign a unique probability that a treated success was
caused by the treatment. An explanation screen that labels a
particular completion as caused by an offer would need additional
information and assumptions.

Assuming monotonic response, $Y(1)\ge Y(0)$ for every customer,
eliminates the harmful response type, forcing $t=2/5$ here.
It implies a benefit fraction $1/5$ and attribution probability
$1/3$. This is a meaningful substantive restriction: assistance
must never prevent a success that would otherwise occur. Delays,
confusion or competing actions can make it implausible in a bank.
The positive average experiment result neither proves nor requires
this restriction.

Expected net value under universal receipt is nevertheless
$100(1/5)-12=8$ USD in every coupling, under the stipulated linear
reward and uniform receipt cost. A sponsor concerned only with
that average reward can compare receipt with no receipt without
resolving $t$. A sponsor imposing a bound on the fraction harmed
cannot. The unresolved individual quantity becomes economically
material through the actual loss function or constraint, rather
than by declaring every nonidentified quantity equally blocking.

\section{Encouragement, receipt and the population are separate targets}
\label{exec17:s3-iv}

Suppose $Z$ is a randomized encouragement, $D$ receipt of assistance
and $Y$ the outcome. Let $D(z)$ be receipt under encouragement
$z$. With stable treatment definitions, assignment independence,
exclusion of a direct encouragement effect, relevance and
monotonicity $D(1)\ge D(0)$, the Wald ratio identifies the mean
receipt effect among compliers. This is the LATE interpretation
developed by Angrist, Imbens and Rubin, rather than a general
population ATE statement.~\citep{angrist1996}

The distinction follows algebraically. Exclusion lets us write
$Y(z,d)=Y(d)$. For binary receipt,
\[
 Y\{D(1)\}-Y\{D(0)\}
   =\{Y(1)-Y(0)\}\{D(1)-D(0)\}.
\]
Monotonicity makes the second factor one for compliers and zero
for never-takers and always-takers. Averaging and using
random assignment gives
\[
 \frac{E(Y\mid Z=1)-E(Y\mid Z=0)}
      {E(D\mid Z=1)-E(D\mid Z=0)}
 =E\{Y(1)-Y(0)\mid D(1)>D(0)\}.
\]
Both the nonzero denominator and the change from receipt to
encouragement are part of the argument. A large fitted first
stage does not verify exclusion or assignment
independence.~\citep{angrist1996}

Consider the authored law with principal strata and observable
response means
\[
\begin{array}{c|c|cc|cc}
 \text{stratum}&P&D(0)&D(1)&E\{Y(0)\}&E\{Y(1)\}\\ \hline
 \text{never}&3/10&0&0&1/5&\text{unknown}\\
 \text{complier}&2/5&0&1&3/10&1/2\\
 \text{always}&3/10&1&1&\text{unknown}&2/5 .
\end{array}
\]
Encouragement is an independent fair bit. Within each stratum
we can construct binary potential responses with the stated
means; product couplings provide one explicit construction.
The receipt means under $Z=0,1$ are $3/10,7/10$;
the outcome means are $3/10,19/50$. The first stage is $2/5$,
the outcome intention-to-treat effect is $2/25$, and LATE is
$1/5$.~\citep{exec17_s3_exact}

What remains unknown is not remedied by calling LATE the effect
for an average customer. Let $n=E\{Y(1)\mid\text{never}\}$
and $b=E\{Y(0)\mid\text{always}\}$. Each may range from zero
to one without changing any observable $(Z,D,Y)$ probability.
The population receipt effect is
\begin{align*}
 \theta_{\mathrm{pop}}
 &=\tfrac3{10}(n-\tfrac15)
   +\tfrac25(\tfrac12-\tfrac3{10})
   +\tfrac3{10}(\tfrac25-b)\\
 &=\tfrac7{50}+\tfrac3{10}(n-b)
   \in[-4/25,11/25].
\end{align*}
Both endpoints and every intermediate value are attainable.
The calculation gives sharp bounds for this particular observed
law under the declared binary IV class, not a claim that this
simple expression is the general IV bounds algorithm.

Universal receipt versus no receipt has net value
in $[-28,32]$ USD per person. Receipt for a complier has mean
incremental net value $8$ USD, but a bank cannot generally
identify each person's principal stratum from the observed
assignment and receipt. It cannot implement treating only
compliers merely because their average effect is
identified.~\citep{angrist1996,exec17_s3_exact}

Changing the feasible action to universal encouragement gives
a different, implementable comparison. Only the $2/5$ compliers
change receipt; with the stipulated success value and receipt
cost, incremental net value is
\[
 100(2/25)-12(2/5)=16/5\ \text{USD}.
\]
A direct encouragement cost of USD 2 reduces this to $6/5$.
Thus a proposal can support the encouragement action under its
assumptions while refusing to claim that universal receipt is
beneficial. The difference is substantive, not merely a choice
between two estimation libraries.

\section{Instrument diagnostics do not supply the missing assumptions}

The IV analysis suggests concrete review questions. Was the
encouragement randomized at a customer or cluster level? Did a
message alter outcomes through attention even when no assistance
was received? Could some customers avoid assistance because of
the encouragement, violating monotonicity? Are nonreceipt and
receipt defined consistently across encouragement groups? These
questions determine the target model before selecting a Wald,
two-stage or structural estimator.

Observed constraints may reject parts of an IV model; a fit that
satisfies them does not certify every unobserved restriction.
Likewise a numerical first-stage diagnostic assesses a feature
of the observable encouragement law. It cannot establish the
population effect among never-takers under receipt, or among
always-takers under nonreceipt. The construction above holds
the entire observed law fixed while changing those two
responses, so no diagnostic based solely on that law can
identify them.~\citep{angrist1996,exec17_s3_exact}

Numerical and inferential weakness also need distinct treatment.
With a fixed positive first-stage denominator $d$, perturbing
the outcome intention-to-treat numerator by $\delta$ changes
the ratio by $\delta/d$. A numerator perturbation of $1/1000$
changes the ratio by $1/400$ at $d=2/5$ and by $1/20$ at
$d=1/50$. At USD 100 per success these are USD $1/4$ and
USD 5 per receipt. This exact algebra explains why small
denominators demand precision. It is not a weak-instrument
confidence procedure: a random denominator and its dependence
on the numerator require an inferential argument of their
own.~\citep{exec17_s3_exact}

A conditional population extrapolation is possible only with
added information or restrictions. A new experiment that can
change receipt in previously unchanged groups supplies a
different information route. A structural effect-homogeneity
assumption supplies a different model route. Repeatedly
estimating the same LATE with greater precision supplies
neither. The proposal should compare feasibility, scientific
credibility and cost of these routes before promising that a
larger historical sample will settle the population question.

\section{Effect sets can already guide a decision}
\label{exec17:s3-bounds}

Return to the first generated law, but now suppose the analyst
receives only $(A,Y)$ and assumes consistency, with no
exchangeability or instrumental restriction. Its observed
cell probabilities are
\[
\begin{array}{c|cc}
 &Y=0&Y=1\\ \hline
 A=0&2/5&1/10\\
 A=1&1/5&3/10 .
\end{array}
\]
The treated successes reveal a mass $3/10$ with $Y(1)=1$.
Among untreated customers, a total mass $1/2$, the treated
potential response is missing. Hence
$E\{Y(1)\}\in[3/10,4/5]$. Similarly,
$E\{Y(0)\}\in[1/10,3/5]$.
The target ATE therefore lies in
\[
 [-3/10,7/10].
\]
For the lower endpoint, give every untreated customer a
missing treated response of zero and every treated customer
a missing untreated response of one. Reverse those choices
for the upper endpoint. Both completions satisfy consistency
and preserve every observed cell; mixtures attain the
intermediate values. Exact enumeration of the sixteen
cellwise extreme completions verifies the construction.
The limits are sharp for this law and model
class.~\citep{exec17_s3_exact}

The point answer $1/10$ obtained with the baseline and
exchangeability assumptions lies inside this set, but the
set does not select it. The two analyses answer questions
under different information and restrictions. More rows
from the unadjusted observational table reduce sampling
error around that table; they do not impose exchangeability.
Recovering a useful baseline is potentially valuable only
if the scientific assignment argument justifies how it
changes the compatible class.

With linear reward, mapping the effect set through
$V=100\theta-12$ gives $[-42,58]$ USD.
For a risk-neutral expected-value decision, a point
probability distribution over compatible models would
produce a mean value, but no such distribution is supplied
by the bounds. Selecting their midpoint adds an arbitrary
decision rule. Under a maximin criterion, no receipt has
value zero and universal receipt has worst value $-42$;
no receipt is preferred in this two-action problem.

Under minimax regret the calculation differs. Relative
to the best of receipt and no receipt in each compatible
world, receipt's worst regret is $42$, whereas no receipt's
worst regret is $58$. If only these deterministic actions
are admissible, minimax regret prefers receipt. Allowing a
randomized receipt fraction $r$, with no interference and
proportional costs, equalizes the endpoint regrets at
\[
 42r=58(1-r),\qquad r=29/50.
\]
The common regret is $609/25$ USD. This is a derived
decision illustration, not an ethical recommendation to
randomize a banking service. Operational, equity and
capacity constraints may exclude that action. It shows
why a causal product must obtain the sponsor's decision
criterion rather than turn an ambiguous effect into an
unexplained binary recommendation.~\citep{exec17_s3_decision}

\section{Identification sets, confidence sets and numerical envelopes}
\label{exec17:s3-numerical}

The preceding ranges describe variation among causal models
with the same exact observable law. They are not confidence
intervals from sampled data. If a separate justified
procedure constructs a confidence region $\mathcal C$
for $P$ under the declared sampling law, an outer effect
region can take the form
\[
 \bigcup_{Q\in\mathcal C}\Theta(Q).
\]
On the event $P\in\mathcal C$, this union contains the true
model's effect, provided the structural class is correct
and the union is evaluated conservatively. That is a
conditional logical implication, not a finite-sample
guarantee for an unspecified estimator. A learned graph
class needs its own retention argument; numerical
approximations to the union need error control as well.

Identification does not by itself control numerical
evaluation. Suppose an identified conditional probability
is the ratio $q=n/d$, with $d>0$. A routine returns
$\widetilde n,\widetilde d$ with genuinely certified
deterministic absolute errors at most $\epsilon$. When
$\widetilde d>\epsilon$ and the probability bounds are
clipped to $[0,1]$, a conservative envelope is
\[
 q\in\left[
 \frac{\max(0,\widetilde n-\epsilon)}{\widetilde d+\epsilon},
 \ \min\left\{1,
 \frac{\widetilde n+\epsilon}{\widetilde d-\epsilon}\right\}
 \right].
\]
The envelope uses monotonicity of a nonnegative ratio.
Correlations between numerator and denominator errors
could tighten it; ignoring those correlations preserves
the outer inclusion. A declared tolerance in an optimizer
is not automatically such a certificate.

Take nominal values
$\widetilde d=10^{-6}$,
$\widetilde n=(3/5)10^{-6}$ and $\epsilon=10^{-7}$.
The envelope is
\[
 q\in[5/11,7/9].
\]
If the comparison risk is exactly $2/5$, the net value
$100(q-2/5)-12$ lies in
$[-72/11,232/9]$ USD. The nominal value is USD 8,
yet the certified numerical envelope crosses zero.
Both nominal masses round to zero at four decimal
places, creating an apparent $0/0$. This is a loss
of numerical information, not proof that the
scientific query is nonidentified.~\citep{exec17_s3_exact}

For comparison, keep the same absolute error but set
$\widetilde d=1/10$ and
$\widetilde n=3/50$. The ratio envelope becomes
\[
 \left[\frac{599999}{1000001},
       \frac{600001}{999999}\right],
\]
with corresponding net values close to USD 8.
The conditioning probability controls how absolute mass
error propagates. In actual data a rare conditioning
event can also cause large sampling uncertainty. The
two effects should be measured separately rather than
combined into an unexplained missing answer.

For the product proposal, this distinction creates
different remedies. A true zero denominator requires
revisiting the target or support assumptions. An
uncertified small positive denominator requires a
numerical diagnostic and a conservative disposition.
Underflow or cancellation may justify rescaling,
log-domain calculation or higher precision followed
by verification. None of these repairs supplies a
missing causal restriction. Conversely, an experiment
that supplies causal information does not repair an
incorrect floating-point implementation.

\section{Choose software by the scientific route it preserves}
\label{exec17:s3-software}

The detailed software comparisons distinguish graphical
adjustment, general intervention identification,
estimation and supplied structural-model simulation.
The synthesis here uses three inspected routes to show
how a product should choose between them; it is not a
claim that these exhaust available identification or
counterfactual software.

DAGitty's inspected adjustment routines reason about a
supplied graph and query. Their useful role is to expose
admissible adjustment choices with the graph's endpoint
meaning intact. They do not learn the graph's scientific
truth from a bank table. The inspected source leaves
some full MAG/PAG validation work as a TODO, and graph
conversion can discard endpoint information. A proposed
integration should therefore validate the admitted
graph class and preserve the original graph alongside
any computational projection. It should reject an
unsupported projection before interpreting an adjustment
result.~\citep{exec13_other_dagitty,exec14_olatent_dagitty}

The inspected y0 high-level API dispatches unconditional
queries to ID and conditioned queries to IDC. Its
default handling turns a particular
\texttt{Unidentifiable} exception into \texttt{None};
strict handling retains that exception. This makes
the exact exception type and query construction
scientifically relevant. A generic missing return
should not erase whether a hedge was obtained, the
input was invalid, an algorithm was unsupported or
an unexpected exception occurred.~\citep{exec13_y0_api,exec13_y0_id,exec13_y0_idc}

The earlier source trace also identifies narrower
concerns. A conditional branch of the inspected query
expression convenience property does not construct
interventions in the same way as the unconditional
branch, although separately stored query fields
remain available. Graph operations can lose an
isolated conditioning node: in the test question
with $A\to Y$ and an isolated $Z$, this may make
computation unavailable even though the stipulated
positive law identifies $P(Y\mid do(A),Z)$.
These are static source findings, not new installed
runtime observations or proofs of mathematical
nonidentification. Section~\ref{exec18:pyid-packages}
adds the separation counterexample and excludes the
selected native conditional path pending repair. A
suitable integration uses explicit query fields,
checks node preservation and admits a joint-first
conditional fallback only with its incompleteness
clearly stated.~\citep{exec13_y0_api,exec13_y0_graph,exec13_y0_dsl,exec13_y0_utils,exec13_y0_rule2,exec13_y0_separation}

The inspected DoWhy model code distinguishes its
IDIdentifier route from its AutoIdentifier route;
choosing an identification method is not equivalent
to choosing an estimator. Its inspected model call
also assigns a local
\texttt{proceed\_when\_unidentifiable} value without
passing that local at the examined call. The proposal
should test the effective option propagation for
the pinned route and preserve an identification
result separately from later estimation. Repeated
refutation exercises on an estimate cannot establish
assumptions that the underlying observation law
does not identify.~\citep{exec13_dowhy}

For these components, adopting a narrow documented
route behind a scientific interface is preferable
to copying an entire package into a new engine
without a parity argument. A direct port changes
the maintenance and verification obligation:
graph surgery, conditional reduction, expression
semantics and failure witnesses all need matched
cases. The existing source-specific rights and
environment reviews remain necessary inputs.
The present evidence supports a prototype comparison,
not a declaration that one stack has been qualified
for every graph or banking regime.

\section{A bounded prototype can resolve implementation choices}
\label{exec17:s3-design}

The six exact law families supply complementary tests
for a proposed public/generated prototype. The
adjustment family checks population weighting and
the distinction between total and controlled direct
effects. The extended bow family checks conditional
identification without an identified joint. The
coupling family checks refusal to turn identified
experimental marginals into identified personal
attribution. The encouragement family checks LATE,
the observable assignment law and the population
effect bounds. The consistency family checks sharp
extreme completions. The rare-event family checks
whether an otherwise valid formula survives numerical
evaluation. Each family has its own expected scientific
answer; agreement on one does not replace the
others.~\citep{exec17_s3_exact,exec17_s3_decision,exec17_s3_design}

Compare components only on a declared shared domain.
An adjustment-only route should be assessed on
adjustment queries; its refusal of a nonadjustment
query should not count as a false identification.
For supported ID/IDC routes, retain the graph,
intervention and condition sets, positivity domain,
formula or witness, and exact expected result.
For a supplied structural-model simulator, retain
the additional mechanisms and exogenous coupling;
agreement with one selected model must not be scored
as identifying the whole observational class.

A meaningful result distinguishes at least a
point functional, a supported effect set, a valid
nonidentification witness, an unsupported scientific
route, invalid input or support, numerical
unavailability, and an execution failure.
These distinctions guide action. Extra assumptions
can shrink a scientific set; a fuller algorithm may
support an unimplemented route; a software repair
may restore a valid answer; additional observations
may reduce sampling uncertainty. Confusing them
would direct the bank toward the wrong investment.

The initial public prototype can use generated
mechanisms with known responses and permissible
public covariates to make the information
structure realistic. Such covariates do not reveal
unobserved potential responses. Before broader
comparisons, the proposal needs a fixed domain,
version and rights decision for each candidate,
matched output semantics and a separate registered
finite-sample study wherever statistical performance
is claimed. Exact identities and deterministic
error envelopes are a necessary reference, not
that study's outcome.~\citep{exec17_s3_design}

Several limits have a known next step without a
known positive answer. Additional primary-source
review can expand software coverage, but uninspected
candidates cannot yet receive an adoption decision.
Generated comparisons can reveal whether the
adapters preserve the mathematical distinction,
but cannot establish a bank's assignment truth.
Domain work and later intervention evidence can
assess plausible instruments, delivery versions,
interference and harm restrictions. Independent
review, accepted loss thresholds, actual precision
requirements and real costs are later inputs.
The proposal can specify those questions and the
evidence that would answer them while proceeding
with the implementation comparisons whose
requirements are already concrete.
\label{exec17:s3-end}


% SOURCE docs/release/2026-10-29/source/docs/writing-repair/2026-10-29/src/87229fb00cac1ea0.tex CHAR 0
\section{The complete observable table constrains an instrument analysis}
\label{exec17:s3-sensitivity}

The instrument example above leaves a practical question unanswered. When the
bank doubts a premise, which conclusions survive, and what additional
information would narrow the answer? A Wald ratio alone cannot answer this.
It discards restrictions in the joint receipt--outcome distribution, while
an assumption change alters the compatible population. We develop the full
binary problem before making either change. Every probability in the following
construction is authored; the sensitivity budgets express hypothetical
scientific judgments rather than facts recovered from bank records.
\citep{exec17_s3_sensitivity}

Let $Z$ be a randomized encouragement, $D(z)$ receipt under encouragement $z$,
and $Y(d)$ success under receipt $d$. For now retain consistency, exclusion
and an encouragement independent of all potential responses, with positive
probabilities for both assignments. Exclusion makes the outcome under the
realized encouragement equal to $Y(D(Z))$. A response type is
$t=(d_0,d_1,y_0,y_1)\in\{0,1\}^4$ with mass $p_t$. The observable
conditional cell probabilities satisfy
\begin{align}
 q_{zdy}
 &=P(D=d,Y=y\mid Z=z) \notag\\
 &=\sum_{t\in\{0,1\}^4}
      p_t\,{\bf1}\{d_z=d,\ y_d=y\}, \qquad
 p_t\geq0,\quad \sum_t p_t=1 .
 \label{exec17:s3-response-map}
\end{align}
The same sixteen masses must generate both assignment groups. Fitting two
unrelated receipt--outcome tables would lose the common-population
restriction supplied by randomization. The assignment probability does not
multiply these conditional cells; it enters the unconditional law and the
sampling precision.

The population receipt effect is the linear objective
$\sum_t p_t(y_1-y_0)$. Its sharp bounds under these assumptions are the
minimum and maximum of this objective subject to
(\ref{exec17:s3-response-map}). This finite optimization represents the
stated information class exactly: every admissible potential-response law
has these masses, and every feasible mass vector constructs an admissible
law by drawing the type independently of $Z$. A constraint on defier mass
is also linear. In contrast, the effect among compliers, the types with
$d_0=0,d_1=1$, is
\[
 \theta_C=
 \frac{\sum_{t:d_0=0,d_1=1}p_t(y_1-y_0)}
      {\sum_{t:d_0=0,d_1=1}p_t}.
\]
This is a linear-fractional target with a scientifically meaningful
denominator. A zero denominator means there is no such stratum; a
numerical convention cannot define its mean effect.

Use the following conditional table, which retains all receipt and outcome
cells of the earlier example:
\[
\begin{array}{c|rrrr}
 &D=0,Y=0&D=0,Y=1&D=1,Y=0&D=1,Y=1\\ \hline
 Z=0&13/25&9/50&9/50&3/25\\
 Z=1&6/25&3/50&19/50&8/25
\end{array}
\]
Each row sums to one. Receipt means are $3/10$ and $7/10$; outcome
means are $3/10$ and $19/50$. Thus the receipt first stage is
$d=2/5$ and the outcome intention-to-treat contrast is
$\tau_Z=2/25$. Their ratio is $1/5$. These two contrasts are summaries
of eight cells, rather than a substitute for them.

\section{Sharp population bounds with monotone receipt}
\label{exec17:s3-monotone-polytope}

Monotonicity, $D(1)\geq D(0)$ for every customer, removes the four
defier response types. Let $N,C,A$ denote never-receivers, compliers and
always-receivers. Their masses are determined by the receipt means:
\[
 \pi_N=3/10,\qquad \pi_C=2/5,\qquad \pi_A=3/10.
\]
Under $Z=1,D=0$ the observed group consists only of never-receivers;
its success mass is $3/50$. Under $Z=0,D=1$ it consists only of
always-receivers; its success mass is $3/25$. Subtracting those
contributions from the other two success cells gives
\[
\begin{array}{c|cc}
 &E\{Y(0)\mid\text{stratum}\}&E\{Y(1)\mid\text{stratum}\}\\ \hline
 N&1/5&n\\
 C&3/10&1/2\\
 A&b&2/5
\end{array}
\qquad 0\leq n,b\leq1 .
\]
Here $n$ and $b$ are conditional means. Below, when principal-stratum
masses vary, we will use different symbols for unconditional success
masses. The distinction prevents multiplication by a stratum weight twice.

The complier effect is $1/2-3/10=1/5$, as the LATE argument predicts
under independence, exclusion, relevance and
monotonicity.~\citep{angrist1996}
The population receipt effect instead is
\begin{align}
 \theta
 &=\tfrac{3}{10}(n-\tfrac15)
   +\tfrac25(\tfrac12-\tfrac3{10})
   +\tfrac3{10}(\tfrac25-b) \notag\\
 &=\tfrac7{50}+\tfrac3{10}(n-b),
 \qquad
 -\tfrac4{25}\leq\theta\leq\tfrac{11}{25}.
 \label{exec17:s3-monotone-ate}
\end{align}
The table supplies neither the outcome of giving receipt to never-receivers
nor the outcome of withholding it from always-receivers.
These are missing counterfactual responses, even with infinite observational
precision and perfect randomization of encouragement.

It remains to prove attainability rather than merely bound a formula.
Within a stratum of mass $m>0$, let $s_0,s_1$ be its success masses
under receipt zero and one, each in $[0,m]$. One admissible coupling is
\begin{equation}
 p_s(y_0,y_1)=
 \frac{a_0(y_0)a_1(y_1)}{m},\qquad
 a_j(1)=s_j,\quad a_j(0)=m-s_j .
 \label{exec17:s3-product-coupling}
\end{equation}
Its four nonnegative cells sum to $m$ and have the required two margins.
Zero-mass strata contribute four zeros. Conditional independence of the two
responses is used here only to construct a witness; it is not imposed on
every model in the bounded class. Other couplings are allowed.

Choose $(n,b)=(0,1)$ for the lower endpoint and $(1,0)$ for the upper
endpoint, and apply (\ref{exec17:s3-product-coupling}) in each of $N,C,A$.
Substituting the resulting twelve possibly nonzero type masses into
(\ref{exec17:s3-response-map}) reproduces every cell of the displayed
table. Convex mixtures of the two type laws reproduce that same table
and attain every intermediate population effect. Equation
(\ref{exec17:s3-monotone-ate}) is consequently a sharp range for this
class. At the stipulated USD 100 per success and USD 12 per receipt,
universal receipt has incremental net range $[-28,32]$ USD per eligible
customer. A positive complier effect therefore does not settle this
population-wide decision.~\citep{exec17_s3_sensitivity}

\section{A bounded departure from monotonicity}
\label{exec17:s3-defiers}

Suppose some customers may refuse the service because the encouragement
was sent. Retain randomization and exclusion, but stipulate
\[
 f=P\{D(0)=1,D(1)=0\}\in[0,1/10].
\]
The observed first stage identifies $\pi_C-f=2/5$, rather than
$\pi_C$ itself. Consequently
\[
 \pi_N=\pi_A=3/10-f,\qquad \pi_C=2/5+f.
\]
The $1/10$ budget is an assumed substantive restriction. Neither the
observed first stage nor an absence of obvious complaints verifies it.
The admitted complier population now varies across compatible models;
it should not be described as a fixed set of known customers.

Let $h_j=P(\text{defier},Y(j)=1)$ be a defier success mass. Subtracting
principal-stratum contributions from all four observed success cells gives
\[
\begin{array}{c|cc}
 &P(\text{stratum},Y(0)=1)&P(\text{stratum},Y(1)=1)\\ \hline
 N&3/50-h_0&u_N\\
 C&3/25+h_0&1/5+h_1\\
 \text{defier}&h_0&h_1\\
 A&u_A&3/25-h_1
\end{array}
\]
where the two unobserved masses satisfy
$0\leq u_N,u_A\leq3/10-f$. Nonnegativity and the stratum sizes imply,
for the declared range of $f$,
\begin{equation}
 0\leq h_0\leq\min(f,3/50),\qquad 0\leq h_1\leq f .
 \label{exec17:s3-defier-feasible}
\end{equation}
There are no further constraints on these margins for this example.
For instance, the never-receiver $Y(0)$ mass is at most its stratum
mass since $3/50\leq3/10-f$, and the always-receiver $Y(1)$ mass
is nonnegative since $h_1\leq1/10<3/25$.
The complier success masses are below $2/5+f$ throughout.
Applying (\ref{exec17:s3-product-coupling}) separately to all four
strata then proves sufficiency. It also reproduces the failure cells,
since the receipt totals and success cells agree.

The complier target becomes
\begin{equation}
 \theta_C(f,h_0,h_1)
 =\frac{2/25+h_1-h_0}{2/5+f}.
 \label{exec17:s3-defier-ratio}
\end{equation}
For fixed $f$ it is increasing in $h_1$ and decreasing in $h_0$.
Its smallest feasible numerator uses $h_1=0$ and
$h_0=\min(f,3/50)$. On $0\leq f\leq3/50$ the lower ratio is
$(2/25-f)/(2/5+f)$, with strictly negative derivative.
On $3/50\leq f\leq1/10$ it is
$(1/50)/(2/5+f)$, again decreasing.
The maximum uses $h_0=0,h_1=f$ and
$(2/25+f)/(2/5+f)$, whose derivative is positive.
Both extrema occur at $f=1/10$, yielding
\begin{equation}
 \boxed{\ \theta_C\in[1/25,9/25]\ },\qquad
 100\theta_C-12\in[-8,24]\ {\rm USD}.
 \label{exec17:s3-defier-sharp}
\end{equation}
The lower witness has $(h_0,h_1)=(3/50,0)$; the upper has
$(h_0,h_1)=(0,1/10)$. Their complier mass is the same, $1/2$.
Product couplings give complete nonnegative response-type laws at both
endpoints; mixtures preserve the table and this denominator, attaining
every intermediate complier effect.

Keeping only $d$ and $\tau_Z$ would give a weaker answer. The contrast
identity under exclusion is
\[
 \tau_Z=(2/5+f)\theta_C-f\theta_{\rm defier}.
\]
Allowing any defier mean effect in $[-1,1]$ gives the outer range
$[-1/25,9/25]$ for the same budget. Its negative lower endpoint
would require $f=1/10$ and $h_0=1/10$, which violates the
observed-cell restriction $h_0\leq3/50$.
Thus fitting a Wald decomposition is insufficient to certify full-law
compatibility. The complete table rules out that negative effect, while
still leaving the USD 12 receipt decision unresolved.

The population effect has a different response to this relaxation:
\[
 \theta=\tfrac7{50}+h_1-h_0+u_N-u_A.
\]
For any admitted $f$, its lower bound is
$-4/25+f-\min(f,3/50)\geq-4/25$ and its upper bound is
$11/25-f+h_1-h_0\leq11/25$.
The monotone endpoint laws remain admitted at $f=0$, so the sharp
population range stays $[-4/25,11/25]$. The larger model class
widens uncertainty for the complier target without widening this
population range. Sensitivity is a property of a specified target and
information class, rather than a single discount applied to every
reported effect.~\citep{exec17_s3_sensitivity}

\section{Exclusion sensitivity needs a different response model}
\label{exec17:s3-exclusion}

A bank message may affect confidence or repayment even when the customer
does not receive the proposed counseling. Exclusion then fails.
The notation $Y(d)$ hides that pathway; use $Y(z,d)$ instead.
In this section restore monotone receipt and retain randomized encouragement,
but allow direct encouragement effects. Define
\[
 \theta_{C0}=E\{Y(0,1)-Y(0,0)\mid C\},\qquad
 \Gamma=E\{Y(1,D(1))-Y(0,D(1))\}.
\]
The first target changes receipt while holding encouragement at zero,
among the receipt compliers. The second averages a direct encouragement
contrast while fixing each customer's receipt to its $Z=1$ value.
It averages over all principal strata. A budget on $\Gamma$ is therefore
not automatically the same as a bound on each individual's direct effect.

Add and subtract $Y(0,D(1))$ inside the randomized contrast:
\begin{align}
 \tau_Z
 &=E\{Y(1,D(1))-Y(0,D(0))\} \notag\\
 &=\Gamma+
   E\{Y(0,D(1))-Y(0,D(0))\} \notag\\
 &=\Gamma+d\theta_{C0}.
 \label{exec17:s3-exclusion-decomposition}
\end{align}
The last equality uses monotonicity to make the receipt-changing group
precisely the compliers. It does not require exclusion.
If the scientific judgment is $|\Gamma|\leq g$, the implication is
\begin{equation}
 \theta_{C0}\in
 \left[\frac{\tau_Z-g}{d},\,\frac{\tau_Z+g}{d}\right]\cap[-1,1].
 \label{exec17:s3-exclusion-outer}
\end{equation}
This is an outer sensitivity envelope. The derivation uses a necessary
contrast equation; it does not prove that every point can satisfy all
eight observed cells in the enlarged $Y(z,d)$ model.
Solving that larger response-type feasibility problem would be required
before calling the envelope sharp.

For the authored table, the implications are
\[
\begin{array}{c|cc}
 g&\theta_{C0}\text{ outer envelope}
   &100\theta_{C0}-12\text{ outer envelope, USD}\\ \hline
 1/50&[3/20,1/4]&[3,13]\\
 1/25&[1/10,3/10]&[-2,18]
\end{array}
\]
Under the smaller budget every compatible receipt effect would exceed
the stipulated cost threshold. Under the larger budget this necessary
envelope cannot settle the decision. It also does not prove that both
signs of net value are attainable in the complete response model.
The budgets are economically consequential assumptions: doubling a
permitted direct contribution changes what this calculation can recommend.
Neither excellent predictive fit nor a small standard error on the
Wald ratio validates those budgets.

\section{A feasible encouragement decision and the next experiment}
\label{exec17:s3-sensitivity-decision}

There is an action whose value uses less causal structure. Suppose the
bank can randomize the encouragement, value every success at USD 100,
pay USD 12 for every receipt, and pay an additional USD 2 when
encouragement is sent. Its incremental expected net from encouragement
one versus zero is
\[
 100\{E(Y\mid Z=1)-E(Y\mid Z=0)\}
 -12\{E(D\mid Z=1)-E(D\mid Z=0)\}-2
 =\tfrac65\ {\rm USD}.
\]
Randomization and consistency for these well-defined assignment policies
identify the two mean contrasts. This calculation remains valid in
the presence of direct encouragement effects or defiers; it does not
divide by the receipt first stage. It concerns offering encouragement,
rather than assigning receipt to unknown compliers. Its positive value
in this exact law is compatible with the unresolved population receipt
range and the sensitive complier comparison.

The information needs now follow from the particular uncertainty.
Recording receipt and outcome jointly preserves the full-cell restrictions
that sharpened the defier calculation. Randomizing encouragement more
precisely estimates those cells but does not observe both receipt
responses for the same person. To narrow population receipt bounds, an
ethically feasible study would need information on responses currently
missing in the never- and always-receiver strata, or defensible additional
restrictions. A design that varies message content separately from service
availability could address the direct pathway, but only if its delivery
and exclusion claims are justified for the new target. Repeating the
same encouragement experiment cannot by itself certify its own exclusion
restriction.

These exact-law results also do not supply sampling confidence intervals.
With estimated cell probabilities, the equality-constrained feasible set
may be empty because of sampling noise even when the population law is
compatible. A proposed inferential extension could form a justified joint
confidence region for the entire conditional table, then optimize over
both that region and admissible response laws. If the region contains
the true table and the assumptions hold, the optimized set contains its
compatible effects. Establishing the region's coverage requires the actual
assignment, clustering and missingness law. Optimizing over eight unrelated
point estimates, or treating a solver's feasibility tolerance as statistical
uncertainty, does not establish that guarantee.

For a small prototype, the full response-type map offers a concrete oracle:
check all cells, distinguish an empty compatible class from a wide effect
set, retain the positive ratio denominator, and return feasible extremal
laws with their bounds. The authored reference evaluates thirty-six
piecewise-domain vertices and separately reconstructs every observable
cell by principal-stratum mixing. Analytic derivatives establish that
the ratio extrema occur among those vertices. This verifies the finite
construction and its attainable ranges; it is a proposed comparison
standard for a future implementation and a basis for choosing a useful
public or generated experiment.~\citep{exec17_s3_sensitivity}
\label{exec17:s3-sensitivity-end}


% SOURCE docs/release/2026-10-29/source/docs/writing-repair/2026-10-29/src/dc5493834897a66f.tex CHAR 0
\section{Identification software must preserve what is and is not identified}
\label{exec18:r-identification}

Suppose a bank is considering an encouragement to make an additional repayment. Let $A$ denote randomized access to a reminder and $Y$ denote repayment by twelve months. The business question may ask for the effect of offering the reminder, whether a particular customer would repay only if reminded, or the range of possible repayment effects when an assumption is unavailable. Those are three different targets. A fitted conditional association, an interventional functional, a cross-world probability and a sharp identified set do not have interchangeable meanings. The proposed interface must name the target and information regime before it calls an R package.

For a supplied directed acyclic graph, \texttt{dagitty} can enumerate sufficient adjustment sets under its documented graphical criteria. That is a useful graph diagnostic; it is not a general algorithm for deriving every interventional functional under latent confounding. The existing \texttt{pcalg} and \texttt{bnlearn} comparison adds equivalence-class and causal-network procedures, but their returned graph or effect list has its own model class and assumptions. In particular, \texttt{pcalg}'s IDA manual assumes a faithful multivariate Gaussian law and a causal DAG without hidden variables, and estimates total effects over admissible DAG extensions. It is neither an arbitrary nonlinear effect region nor a confidence set for a bank decision. The \texttt{ivmodel} confidence-set interface answers another question again: its documented AR/CLR inversion returns interval components under its declared IV design, with CLR described as robust to weak instruments. The proposed wrapper must preserve these set components and admit infinite endpoints where the inversion yields them. A single label such as “causal effect” would erase these distinctions~\citep{exec13_other_dagitty,exec13_pcalg,exec13_bnlearn,v12_sw_ivmodel}.

The R package \texttt{causaleffect} supplies symbolic ID/IDC functionality and selected transportability and selection-recovery procedures. Its pinned R function accepts outcome, treatment, optional conditioning variables and a graph; it can return an expression or a derivation object. The selected implementation has an important adapter consequence: with its default \texttt{stop\_on\_nonid=TRUE}, nonidentification raises an error; with the option false, a failed identification can be represented by \texttt{id=FALSE} in the step result, an empty string, or \texttt{NULL}, depending on the other output options. A production wrapper must distinguish that documented algorithmic result from parse errors, unsupported graphs and runtime exceptions. It should also retain the graph class and functional, not just a Boolean or formatted string~\citep{exec18_rid_causaleffect}.

The pinned \texttt{dosearch} package asks whether an interventional target follows from multiple observational or experimental distributions and a graph, including selected transport, selection and missingness mechanisms. Its 2021 paper describes a sound search over do-calculus rules, with completeness tied to the specified problem class. Its Table 1 distinguishes complete ID/IDC cases, unknown completeness for general multiple-source identification, and an incomplete missing-data extension. The package documentation expressly says its missing-data search is not complete. More subtly, the pinned C++ source applies a positive time limit by returning early from the search; its initializer then decides identifiability from whether a target distribution was found and does not expose a separate timeout state in that result. Therefore, under that code path, a limit-truncated search can be represented like “not identifiable by this search.” An adapter must not promote that output to a proof of nonidentifiability. It must either disable the limit for the bounded admitted problem, instrument a distinct timeout/unknown state, or report the uncompleted search as unknown. This is a source-level review of a pinned revision, not an executed package trial~\citep{exec18_rid_dosearch,exec18_rid_dosearch_paper}.

Counterfactual identification is a further level. \texttt{cfid} implements ID*/IDC* for counterfactual conjunctions by constructing a parallel-worlds graph with shared latent variables and then applying recursive identification. Its function defaults to identifying from interventions, not from observations; the observational level must be requested or obtained through the documented alternatives. A true \texttt{id} field is not sufficient on its own: an inconsistent conjunction may be identifiable with probability zero, and a conditional query may be undefined when its conditioning event has zero probability. An adapter therefore retains the requested data level, identifying formula, returned \texttt{undefined} field and query itself; an inconsistency represented by a zero formula must retain that meaning. The package paper explains why shared latent variables connect the counterfactual worlds; treating them as independent simulated populations would change the query~\citep{exec18_rid_cfid,exec18_rid_cfid_paper}.

When point identification fails, bounds may still be useful, but their scope is not universal. \texttt{causaloptim} encodes certain graph-defined response-function problems as linear programs and permits added constraints; its documented graph attributes and admissible left/right structure delimit the class. \texttt{bpbounds} implements binary treatment and outcome bounds for an instrument with two or three categories, in a single trivariate table or a specified two-source bivariate design, with separate monotonicity outputs. Neither manual claims to deliver a sampling confidence region for every reported plug-in endpoint. The distinction matters: an identified set describes values compatible with a population law and assumptions; sampling uncertainty concerns how an estimator or set varies across repeated samples. A posterior interval from a separate Bayesian model is different again. The software contract should return these objects with different types and labels, including empty, unbounded, undefined and unknown states where relevant~\citep{exec18_rid_causaloptim,exec18_rid_bpbounds}.

\texttt{CausalQueries} is a useful adjacent tool rather than a substitute for a graphical identification test. Its documentation defines binary causal models, updates model parameters and supports queries over priors, parameters or posterior draws. Its query manual calls the default interval a “confidence interval” under the statistics argument, but names the interval endpoints as credible summaries and allows a prior or posterior distribution as the input. The proposed interface must attach the distribution and uncertainty meaning explicitly instead of inheriting that ambiguous label. Such a posterior answers a question conditional on the declared graph, response-type parameterization, restrictions, prior and update model. It need not equal a likelihood-identified set, and a narrow credible interval cannot certify that omitted graphs or cross-world dependence have been resolved. For the proposed library, the package is a reference for probabilistic model-query syntax and an optional R adapter candidate. The described package and its compiled Stan dependencies have not been installed, executed or cleared for redistribution in this review~\citep{exec18_rid_causalqueries}.

\subsection{An exact repayment example separates four outputs}

Let $Y_0,Y_1\in\{0,1\}$ be repayment under no reminder and reminder. Randomization identifies $\Pr(Y_0=1)=1/4$ and $\Pr(Y_1=1)=3/4$ in this deliberately stipulated population example. Write $q_{ij}=\Pr(Y_0=i,Y_1=j)$ and let $t=q_{11}$. The observed arm risks imply
\begin{align}
 q_{11}&=t, &q_{10}&=\frac14-t,\notag\\
 q_{01}&=\frac34-t, &q_{00}&=t,\qquad 0\leq t\leq\frac14. \label{exec18:r-fiber}
\end{align}
The intervention effect is point identified throughout this fiber:
\begin{equation}
 \operatorname{ATE}=\mathbb E(Y_1-Y_0)=\frac34-\frac14=\frac12.
 \label{exec18:r-ate}
\end{equation}
But the retrospective “necessity” fraction among would-be responders under the reminder is
\begin{equation}
 \operatorname{PN}=\Pr(Y_0=0\mid Y_1=1)
 =\frac{q_{01}}{q_{01}+q_{11}}=1-\frac{4t}{3}\in\left[\frac23,1\right].
 \label{exec18:r-pn}
\end{equation}
The average effect is a point, while the counterfactual fraction is only set identified. Assuming monotonicity $Y_1\geq Y_0$ removes the response type $q_{10}$, forces $t=1/4$ and narrows PN to $2/3$. That is an added scientific restriction, not information generated by a software call.

Every value in this interval is attainable: choose the corresponding $t$ and sample the four response types with probabilities $q_{ij}$, independently of randomized assignment. This supplies a constructive sharpness argument. The population risks above are stipulated constants; observing one success versus three in two samples of four would not establish those risks without sampling uncertainty.

The same distinction survives Bayesian updating. Consider the three compatible worlds $t\in\{0,1/8,1/4\}$. With four randomized units in each arm and the stated one-of-four versus three-of-four successes, the likelihood of the ordered outcomes is $(1/4)^2(3/4)^6=729/65536$ in every world. Thus the data do not change their relative weights. A uniform prior on these three worlds gives posterior mean PN $5/6$; prior weights $(1/6,1/3,1/2)$ give $7/9$. These are exact project calculations, not output from \texttt{CausalQueries} and not a defect in Bayesian computation: the difference records prior sensitivity along a likelihood-flat fiber. The aggregate binomial likelihood is $729/4096$ in every world, including the common combinatorial factor.

For an illustrative retrospective accounting rule, assign USD 1,000 to a responder whose repayment is attributed to the reminder and subtract USD 800 per responder. The compatible net amount per responder ranges from $1000(2/3)-800=-\$400/3$ to $1000(1)-800=\$200$. The two posterior means are USD $1000(5/6)-800=100/3$ and USD $1000(7/9)-800=-200/9$, respectively. These stipulated amounts are not observed economics or ROI, and attribution after the outcome must not be misread as prospective treatment selection on future repayment. The executable Fraction reference enumerates all three response-type distributions and checks every displayed likelihood and monetary value~\citep{exec18_rid_exact}.

\subsection{Conditional reuse and the boundary of a negative answer}

The lending case makes the graph restriction concrete. Suppose unrecorded risk $U$ affects both the offer $A$ and repayment $Y$. In a deliberately simple example, take $U\sim\operatorname{Bernoulli}(1/2)$ and $A=U$. Two structural models, $Y=A$ and $Y=U$, give the same observed $A=Y$ distribution, yet their average effects of intervening on $A$ are respectively one and zero. The graph allowing $U\to A$, $U\to Y$ and $A\to Y$ contains both models, because an allowed parent may have zero effect. No precision of fitting that observed distribution selects the intervention answer. This is an identification failure relative to that information and graph class. The earlier randomized example changes the information and identifies the average effect, while leaving cross-world attribution ambiguous.

For this use, the proposed adapter result retains three epistemic states: identified with a functional or point; not identified by a complete algorithm for its declared graph/data class; and unknown because the class is unsupported or computation ended before a certificate. Invalid input is a separate validation failure, with no scientific conclusion. Sharp identified sets, sampling confidence sets and Bayesian posterior summaries have separately named result types and assumptions. A zero-probability conditioning event is an undefined query, not a negative identification certificate.

The practical proposal is to retain \texttt{causaleffect} as a candidate symbolic R service and \texttt{cfid} as a counterfactual-query service, preserving their query objects and derivations. \texttt{dosearch} complements those services for multiple source distributions, but an adapter must record admitted search rules and termination before interpreting a negative result. Disabling a time limit does not make an incomplete rule system complete. Porting any of these algorithms requires an independently checked mathematical contract and differential fixtures; matching a formatted expression alone would miss a changed conditioning event or data level. For that reason direct source ports are deferred while optional process-bound R adapters are designed. This is a proposed interface boundary, not an executed integration or a conclusion about license compatibility~\citep{exec18_rid_causaleffect,exec18_rid_cfid,exec18_rid_dosearch}.

The restricted bounds engines offer a different reuse choice. \texttt{causaloptim} is a candidate optimization reference when its response-function class is appropriate; \texttt{bpbounds} supplies a smaller binary-IV reference with explicit monotonicity and two-source assumptions. The proposal rejects either as a universal sharp-bounds backend. It retains \texttt{CausalQueries} as an optional model-query path when priors and sensitivity are substantively justified, while rejecting a posterior interval as a replacement for identification. Dependency and license observations inform these choices but do not qualify redistribution, linked dependencies, numerical behavior or installed package execution~\citep{exec18_rid_causaloptim,exec18_rid_bpbounds,exec18_rid_causalqueries}.

These comparisons establish distinct candidate roles and failure semantics. Coverage of the wider R ecosystem remains incomplete; selected packages and their inspected versions do not exhaust the available methods. The exact examples establish the stated mathematical distinctions and illustrative accounting consequences, while implementation and broader source coverage remain future work.


% SOURCE docs/release/2026-10-29/source/docs/writing-repair/2026-10-29/src/25c488b856f3d0d8.tex CHAR 0
\section{Preserve the adjustment question across the DAGitty wrapper}
\label{exec18:gp-graph}

A bank analyst who supplies a causal diagram still needs to specify which
intervention and outcome the diagram is meant to relate. A graph editor, a
test of an adjustment set, an enumeration of such sets and an estimator of an
effect answer different questions. Our proposed use of DAGitty is deliberately
narrow: a supplied observed directed acyclic graph, one exposure, one outcome
and a total-effect adjustment query. The proposed first implementation would
call a pinned original-language R wrapper and preserve its graph and query
semantics. This is a proposal selection, not an installed adapter or a claim
that the bank's supplied arrows are true.
\citep{exec18_gp_dagitty}

The selected repository commit is dated 6 December 2024; its R description
declares version 0.3-4 and date 7 December 2023. Those dates identify the
inspected historical source and do not establish present support. The README
describes the shared JavaScript core and the R/V8 boundary. The R adjustment
entry points check that exposure and outcome variables are supplied and stop
with an error if either is absent. The inspected JavaScript Boolean check
instead returns false in that case. An application that interprets every false
as a scientific finding that adjustment is impossible would turn a missing
question into a causal conclusion. The adapter must validate the question
before interpreting the upstream value. \citep{exec18_gp_dagitty}

For a complete supported query, three distinct outputs remain necessary.
A valid empty covariate set means no adjustment is required; an empty
collection of valid sets means that this adjustment search supplied none;
an incomplete search or invalid request remains unresolved. In particular,
\(\{\varnothing\}\) and \(\varnothing\) have different meanings. The first
contains one set and the second contains none. Even an exhaustive failure
to find an adjustment set does not exclude identification through a method
outside covariate adjustment. The documented \texttt{max.results} argument
caps minimal-set enumeration. A capped listing cannot be reported as a list
of all valid choices, and an interrupted or unqualified computation cannot
be promoted to a proof of nonidentification. These distinctions belong in
the returned result and its completeness field. \citep{exec18_gp_dagitty}

\subsection{Three banking laws expose different conversion errors}
\label{exec18:gp-graph-laws}

Consider three separate project-authored structural laws. Their probabilities
are stipulated; they are mathematical references for a future adapter, not
observed customers or executed DAGitty examples. In the first, \(A\) is a
binary intervention, \(Y\) a binary favorable repayment outcome and \(Z\) a
baseline readiness attribute. The graph is \(Z\to A\), \(Z\to Y\),
\(A\to Y\), with mutually independent structural randomizers and
\begin{equation}
\label{exec18:gp-confounding}
 \begin{aligned}
 \Pr(Z=1)&=\frac12, &
 \Pr(A=1\mid Z=z)&=\frac{1+2z}{4},\\
 \Pr(Y=1\mid A=a,Z=z)&=\frac{1+2a+z}{8}.
 \end{aligned}
\end{equation}
Every relevant treatment probability is positive. Under this declared graph
and law, standardizing over the population distribution of \(Z\) gives
\(\Pr\{Y(0)=1\}=3/16\) and \(\Pr\{Y(1)=1\}=7/16\). Their difference is
\(1/4\). The unadjusted observed risks are instead \(5/32\) and \(15/32\),
whose difference is \(5/16\). Here the supplied baseline variable is needed
for the adjustment interpretation. A wrapper that drops the selected set
would change the target calculation. Exact enumeration of all eight joint
cells reproduces both pairs. \citep{exec18_gp_graph_reference}

The second law shows why adding variables is not a universal remedy. Let
\(A\) and an unobserved resilience variable \(U\) be independent
Bernoulli\((1/2)\) variables, let \(Y=U\), and let documentation completion
\(C\) satisfy
\begin{equation}
\label{exec18:gp-collider}
 \Pr(C=1\mid A=a,U=u)=\frac{1+2a+2u}{8},
 \qquad A\to C\leftarrow U\to Y.
\end{equation}
There is no directed path from \(A\) to \(Y\), so the total effect is zero.
The empty adjustment set is valid. Nevertheless, the observed repayment
contrasts within \(C=0\) and \(C=1\) are respectively \(-1/24\) and
\(-1/8\). Standardizing these contrasts over the marginal distribution
of \(C\) gives \(-7/96\). All eight \((A,U,C)\) cells have positive probability;
the discrepancy is caused by conditioning on the collider, rather than by
a missing-support cell. A business workflow that automatically adds a
post-intervention documentation field can therefore introduce a spurious
adverse relationship. The graph and the information clock must survive
conversion together. \citep{exec18_gp_graph_reference}

In the third law, \(A\) is randomized, \(M\) is subsequent engagement,
and the graph contains \(A\to M\to Y\) and \(A\to Y\). Set
\(\Pr(M=1\mid A=a)=(1+2a)/4\) and
\(\Pr(Y=1\mid A=a,M=m)=(1+a+2m)/8\). The total-effect contrast is
\(1/4\). Holding the marginal distribution of \(M\) fixed while
standardizing the conditional contrasts gives \(1/8\). Thus a descendant
can change a total-effect calculation even when positivity holds. This
second contrast is a different functional of the stipulated law; it is
not a substitute for the total effect simply because both are numbers.
The wrapper's documented total/direct-effect switch must be an explicit
part of the request. A new claim about a direct-effect estimand would
need its own assumptions and interpretation. \citep{exec18_gp_graph_reference,exec18_gp_dagitty}

\subsection{Choose a runtime boundary without erasing graph uncertainty}

The wrapper documents graph classes beyond observed DAGs, but the proposed
first admission does not inherit every advertised class. The inspected
JavaScript validation code leaves proper MAG/PAG validation as a TODO, and
the inspected \texttt{pagToPdag} transformation changes endpoint marks and
graph type. Replacing a partial ancestral graph by a directed representation
may discard the ambiguity relevant to an intervention. A future extension
must preserve endpoint marks, latent and selection meanings, class validity
and the scope of the adjustment criterion. Syntax acceptance alone does not
supply those guarantees. Restricting the initial adapter to an observed DAG
keeps this unresolved composition visible. \citep{exec18_gp_dagitty}

Original-language reuse is preferred to an immediate port because it gives
a concrete comparison surface: parse the same supplied graph, validate the
same query, and compare typed adjustment results against independent finite
references. This is an engineering rationale, not evidence that a wrapper
already conforms. A direct JavaScript integration remains an alternative
when its preconditions and packaging are independently qualified; a port
would additionally need algorithmic equivalence, graph-class preservation
and a maintenance plan. Browser editing can support human review, but it
does not establish that a discovered or expert-supplied diagram is causal.
No browser, R, Node or upstream test execution supports the current choice.
\citep{exec18_gp_dagitty}

Rights also follow the selected components. The root contains GPL version 2
text; the R description declares GPL-2; one inspected helper has a
GPL version 2-or-later header, and other helpers declare MIT terms. The tree
contains vendored Underscore versions 1.8.3 and 1.13.1 and a separate test
dependency range. Those observations do not establish which copy a proposed
runtime loads, a vulnerability, or a complete transitive-rights inventory.
The release artifact, dependency resolution, notices, intended distribution
and source obligations must be examined together. Neither a separate process
nor a public download alone settles the rights of a particular arrangement.
The defensible proposal is pinned reference reuse with executable admission
deferred until that bounded trial and rights review are complete.
\citep{exec18_gp_dagitty}


% SOURCE docs/release/2026-10-29/source/docs/writing-repair/2026-10-29/src/6745e9f384757b9a.tex CHAR 0
\section{Python identification: the question determines the reusable capability}
\label{exec18:pyid-question}

A lending analyst may ask whether changing an offer improves average repayment,
or whether a particular customer who failed to repay would have succeeded under
another offer. These questions share variables but require different information.
Let $A$ be a binary intervention and $Y(a)$ the repayment outcome under its
assignment to $a$. An average contrast $E\{Y(1)-Y(0)\}$ uses two marginal
intervention distributions. The retrospective question
$P\{Y(1)=1\mid A=0,Y=0\}$ also links the factual outcome to that person's
unobserved outcome under another assignment. A package capable of computing the
first quantity does not thereby identify the second. Nor does a representation
of a counterfactual variable establish an algorithm for identifying its law.
The distinction changes the proposed product boundary: a request must state its
query class and available information before an engine is selected.
\citep{exec13_y0_dsl,exec18_pyid_contract}

\subsection{Four complementary source boundaries}
\label{exec18:pyid-packages}

The inspected DoWhy revision separates model specification, identification and
estimation. Its identification method chooses between an ID identifier and an
automatic identifier. It can also construct a graph from supplied inputs, so
the proposed banking route requires an explicit reviewed graph instead of
treating that convenience as elicitation evidence. In the selected revision,
the function assigns the \texttt{proceed\_when\_unidentifiable} option but does
not forward it to either identifier; validation must inspect the returned
estimand and assumptions. This is a static source finding, not a demonstrated
runtime failure. Optional reuse of the pinned identification boundary remains
the preferred initial path. Section~\ref{sec:exec24-gcm-reuse} supplies the
separate GCM source review, including its structural assumptions and
factual-noise reconstruction. Its narrower reference-only decision does not
inherit qualification from this identification interface.
\citep{exec13_dowhy}

The selected pgmpy adjustment interface offers a narrower complementary role:
generation of candidate sets for a specified exposure and outcome. Candidate
generation must be followed by independent causal validation. The reviewed
implementation checks exposure parents in a way that can accept a supplied
mediator set when the exposure has no parents; its \texttt{minimal\_variance}
branch is advertised but raises \texttt{NotImplementedError}. An adapter must
exclude that branch and validate the chosen total-effect adjustment criterion.
Neither a minimal separator nor a successful role check guarantees a globally
minimum-variance adjustment set. These restrictions do not establish anything
about uninspected pgmpy modules.~\citep{exec13_pgmpy_reuse_primary}

The pinned y0 implementation gives a symbolic ADMG identification route.
Its conditional algorithm invokes a rule-2 exchange test and otherwise
normalizes an identified joint distribution. The inspected separation path has
the previously developed bidirected-collider counterexample: conditioning on
the collider can give probability $9/10$, whereas replacing the relevant
variable by intervention gives $1/2$. The proposed native conditional path
therefore remains disabled. Identifying the joint distribution and then
normalizing is a possible restricted alternative when that joint is identified
and its denominator is positive. This sufficient route is not a completeness
claim for conditional identification: failure to identify the entire joint
does not prove failure to identify a particular conditional target.
\citep{exec13_y0_id,exec13_y0_idc,exec13_y0_rule2,exec13_y0_separation,exec13_y0_oracles}

Ananke adds a differently organized source reference. The selected authentic
\texttt{ananke-causal} 0.5.0 distribution was uploaded on 5 December 2023; it is
not presented as the newest release. Its \texttt{OneLineID} object fixes
intervention vertices, finds relevant outcome ancestors and checks whether
each district's kernel can be obtained by fixing other vertices. The \texttt{id()}
method returns a Boolean. The \texttt{functional()} method returns a string
describing the fixing expression and prints fixing-order diagnostics. That
string is not yet a typed portable expression, a fitted estimator or an
uncertainty statement. Proposed reuse would preserve variable identities,
sums, products and fixing operations in a restricted expression representation;
it would never evaluate arbitrary returned text as program code.
\citep{exec18_pyid_one_line,exec18_pyid_contract}

Ananke's experimental identification variants also show why exception names
cannot be the scientific result contract. \texttt{OneLineAID} raises
\texttt{NotIdentifiedError} when the experimental graphs fail its ancestral
precondition; \texttt{OneLineGID} uses the same exception for nonconforming
experimental graphs. A precondition failure must be distinguished from a
valid query proved nonidentified in a declared model class. The source's
supplied-model witness routines compare observational marginals and
interventional distributions of two supplied models. They do not establish
automatic witness construction for every query or arbitrary cross-world
identification.~\citep{exec18_pyid_one_line}

Other Ananke components answer different questions. Its oracle intervenes in a
supplied true model and explicitly states that identification is not considered.
Conditional normalization inside that oracle cannot be cited as identification
from observed data. The missing-data routine addresses full-law identification
with explicit \texttt{R\_} and \texttt{X\_} naming and colluding-path checks.
The \texttt{CausalEffect} estimation interface targets the average contrast
$E\{Y(1)-Y(0)\}$ and selects among adjustment-, primal- and nested-fixability
strategies. A file named \texttt{counterfactual\_mean.py} is therefore evidence
about that average-effect interface, not a guarantee for a customer's joint
potential outcomes.~\citep{exec18_pyid_oracle,exec18_pyid_missing,exec18_pyid_mean}

\subsection{What even a randomized study leaves undetermined}
\label{exec18:pyid-coupling}

Consider an exact illustrative population, rather than a fitted bank dataset.
Let $A$ and a latent response type $U$ be independent fair binary variables,
and impose consistency $Y=Y(A)$ and no interference. In the stable model
$Y(0)=U$ and $Y(1)=U$. In the switching model $Y(0)=U$ and $Y(1)=1-U$.
Both are compatible with an allowed $A\to Y$ causal model when faithfulness
is not imposed. The first has no individual response to assignment; in the
second, improvements and harms cancel in the mean. Direct enumeration gives
\begin{equation}
\begin{aligned}
 P(A=a,Y=y)&=\tfrac14 &&(a,y\in\{0,1\}),\\
 P\{Y(0)=1\}&=P\{Y(1)=1\}=\tfrac12,\\
 E\{Y(1)-Y(0)\}&=0.
\end{aligned}
\label{exec18:pyid-equal-laws}
\end{equation}
Thus an arbitrarily large randomized study can recover the marginal laws in
this example without identifying their within-person coupling. Among factual
untreated failures, $A=0,Y=0$ implies $U=0$. The stable model then has
$Y(1)=0$, whereas the switching model has $Y(1)=1$. Consequently,
\begin{equation}
 P\{Y(1)=1\mid A=0,Y=0\}=
 \begin{cases}0,&\text{stable model},\\1,&\text{switching model}.\end{cases}
\label{exec18:pyid-conditional}
\end{equation}
Both conditioning events have probability $1/4$, so this disagreement is not
created by conditioning on an impossible event. Exact rational enumeration
retains the complete observed and response-type tables for both models.
\citep{exec18_pyid_coupling}

The ambiguity can be described sharply. Write $p_{ij}=P\{Y(0)=i,Y(1)=j\}$
and let $t=p_{01}$ be the proportion helped by the intervention. The two
marginal constraints and normalization force
\begin{equation}
 (p_{00},p_{01},p_{10},p_{11})
   =\left(\tfrac12-t,t,t,\tfrac12-t\right),
 \qquad 0\leq t\leq\tfrac12.
\label{exec18:pyid-sharp}
\end{equation}
Nonnegativity gives the necessary interval. Conversely, for every $t$ in that
interval, take a latent variable with these four response-type probabilities
and randomize $A$ independently. The resulting structural model attains the
same observed law and intervention marginals. This construction proves
attainability of every interior point as well as both endpoints. Therefore
$P\{Y(1)=1,Y(0)=0\}$ has sharp bounds $[0,1/2]$, and randomization yields
\begin{equation}
 P\{Y(1)=1\mid A=0,Y=0\}
 =\frac{p_{01}}{P\{Y(0)=0\}}=2t\in[0,1].
\label{exec18:pyid-bounds}
\end{equation}
This is an identified set under a stated information model, not a confidence
interval. Sampling uncertainty would require a further procedure that accounts
for estimation of the marginal restrictions.~\citep{exec18_pyid_coupling}

Suppose, solely to interpret the target, a successful recovery is worth
USD 400 and changing the original intervention costs USD 50. The retrospective
conditional net amount is $400(2t)-50$, with sharp range from minus USD 50 to
USD 350. The two models reverse the sign of this accounting despite identical
average causal effects. This quantity does not establish an implementable
prospective policy: the factual failure used to select the group was not
known before the original assignment. A future servicing decision would need
its own time-indexed action, available history and outcome horizon. The
example motivates a refusal or an explicit identified set, not hindsight-based
deployment.~\citep{exec18_pyid_coupling}

Additional scientific restrictions can change the answer. Monotonicity
$Y(1)\geq Y(0)$ rules out the harmed response type, so $p_{10}=t=0$ under
these equal marginals. It then forces the stable coupling. Randomization alone
does not supply that restriction. Repeated observations, mechanistic evidence
or a richer structural model would have to justify their own connection across
worlds and time; simply collecting more independent assignments does not reveal
both incompatible assignments for the same occasion. The proposed service
must record the added assumption and its scientific owner before narrowing
the set.~\citep{exec18_pyid_coupling,exec18_pyid_contract}

\subsection{Integration choices follow semantics and dependency constraints}
\label{exec18:pyid-integration}

The reuse proposal is deliberately component-specific. DoWhy supplies a
candidate explicit-graph identification boundary. pgmpy supplies candidate
adjustment sets with independent admissibility checks. y0 supplies a candidate
symbolic ID route, with its native conditional path excluded pending a
separation repair. Ananke supplies a district-fixing reference and a possible
isolated comparator. Directly adopting all four packages would not create one
uniform counterfactual service. Porting their source would transfer symbolic,
graph and maintenance obligations to this project without a demonstrated
deployment advantage; it is not selected initially.
\citep{exec13_dowhy,exec13_pgmpy_reuse_primary,exec13_y0_id,exec18_pyid_one_line,exec18_pyid_contract}

There is also a concrete environment conflict. The selected Ananke release
declares NumPy at least 1.23 and below 2, statsmodels at least 0.13.2 and below 0.14,
and pgmpy at least 0.1.21 and below 0.2. The selected DoWhy revision requires
NumPy above 2 and statsmodels at least 0.14; the selected pgmpy revision declares
NumPy at least 2 and statsmodels at least 0.14.5. These declared intervals cannot
all be satisfied in one unchanged environment. Separate pinned processes
communicating through a named-variable result contract are a design option,
not a tested environment resolution. Ananke also imports Sympy in the
identification source without explicitly declaring it in the inspected wheel
metadata; whether a resolved transitive dependency supplies it remains to be
checked. No import failure has been observed here.
\citep{exec18_pyid_metadata,exec18_pyid_one_line,exec13_dowhy_project,exec13_pgmpy_reuse_primary}

The proposed boundary validates cycles, node identity, edge type, latent status
and experimental preconditions before dispatch. Ananke's graph constructor
contains segregation and cycle checks, but their presence and the nearby
correctness TODO do not constitute either a qualification or a demonstrated
defect. Engine diagnostics are retained separately from scientific disposition.
An identified expression, an identified set, a model-specific nonidentification
result, an unsupported question, invalid input and an engine failure remain
distinct outcomes. An empty valid adjustment set must not become a failed query,
and an exception must not become a proved impossibility.
\citep{exec18_pyid_graph,exec18_pyid_one_line,exec18_pyid_contract}

Rights are similarly tied to the actual component and distribution. The
selected DoWhy source carries MIT text and y0 carries BSD three-clause text.
The Ananke archive contains two byte-identical Apache 2.0 license files, while
its metadata also includes an Other/Proprietary classifier. The conflicting
declarations remain visible for distribution review rather than being silently
resolved by a package badge. Notice retention, modified-file marking where
applicable and dependency-specific obligations belong in the selected
distribution design. Process separation does not settle those questions.
\citep{exec13_dowhy_license,exec13_y0_license,exec18_pyid_license,exec18_pyid_metadata}

This comparison leaves explicit survey work. The selected DoWhy GCM paths
receive their own source review in Section~\ref{sec:exec24-gcm-reuse}; that
review does not complete native qualification or the wider census.
Implementations of general counterfactual identification and Python
partial-identification optimizers still need authenticated technical
sources, model classes and output checks. A solver would need to distinguish
a feasible witness from a certified sharp endpoint; a counterfactual engine
would need to state what structural information permits abduction and coupling.
The four inspected families support the bounded reuse choices described here.
They do not establish exhaustive Python coverage, installed adapters or
qualified banking decisions.~\citep{exec18_pyid_contract}


% SOURCE docs/release/2026-10-29/source/docs/writing-repair/2026-10-29/src/8ac3faa3165e4001.tex CHAR 0
\section{Distinguishing identified effects from unresolved comparisons in Java}
\label{exec19:javaid-scope}

A lending analyst may ask whether approving a loan changes repayment, or whether
a particular declined borrower would have repaid if approved. Those questions
differ from updating a supplied Bayesian network after observing borrower
attributes. They also differ from collecting regression coefficients over
possible graph orientations. A Java integration must preserve these distinctions
before it exposes a probability or an effect. Interface names alone cannot
determine the scientific contract.

The selected Tetrad source, version 7.6.11-SNAPSHOT at preserved revision
\texttt{b89afaf952a4}, contains an \texttt{Identifiability} interface in addition
to \texttt{Ida} and \texttt{PdagPagIda}. The identification implementation
constructs c-components and Q-expressions and evaluates them using an
observational joint-probability representation. It is a distinct candidate
from continuous-data effect enumeration. The source header names a 2002
algorithm report, but the report's proof has not been inspected for this
comparison. Code observations establish interface and failure-path requirements;
they do not establish theorem completeness or unrestricted latent-variable
identification. Nine selected original files support a bounded Tetrad comparison,
not an exhaustive census of Java software.~\citep{exec19_javaid_source}

\subsection{Observation, intervention and individual counterfactuals}
\label{exec19:javaid-witnesses}

Consider two authored lending mechanisms. Let $U$ be a binary borrower state,
$A$ approval, and $Y$ a binary repayment outcome. Let $U$ have probability
$1/2$ of being one, and let independent noise variables $N,E$ each have
probability $1/4$ of being one. Write $\oplus$ for exclusive-or and set
\begin{equation}
 A=U\oplus N,\qquad
 \mathcal M_1:\ Y=U,\qquad
 \mathcal M_2:\ Y=A\oplus E.
 \label{exec19:javaid-models}
\end{equation}
Both mechanisms are compatible with the allowed supergraph
$U\to A$, $U\to Y$, $A\to Y$ when an allowed arrow may have zero functional
effect. This construction does not impose faithfulness. Both yield the same
strictly positive observed joint law,
\begin{equation}
 \begin{array}{c|cc}
  &Y=0&Y=1\\\hline
 A=0&3/8&1/8\\
 A=1&1/8&3/8
 \end{array},
 \qquad P(Y=1\mid A=1)=3/4.
 \label{exec19:javaid-law}
\end{equation}
Nevertheless, the intervention risks are $(1/2,1/2)$ in $\mathcal M_1$ and
$(1/4,3/4)$ in $\mathcal M_2$, ordered by approval zero and one. Their average
causal effects are zero and $1/2$. The positive observed law does not select
one of these effects under the declared model class. These are two explicit
witnesses to nonuniqueness, not an enumeration of a sharp identified set. A
conditional probability cannot be relabeled as an intervention risk merely
because both are between zero and one.~\citep{exec19_javaid_reference}

For a purely illustrative monetary conversion, stipulate USD 200 per unit
increase in repayment probability and an approval-action cost of USD 20. The
two mechanisms give incremental net values
\begin{equation}
 200\{E[Y\mid\mathrm{do}(A=1)]-E[Y\mid\mathrm{do}(A=0)]\}-20
 \in\{-20,80\}\quad\text{USD per case}.
 \label{exec19:javaid-money}
\end{equation}
The sign changes without changing the observed joint law. These amounts are
stipulated arithmetic, not loan-pricing evidence, a measured bank return or
the valuation of a real approval policy. Additional causal restrictions or
intervention evidence would be needed to choose a mechanism for decision use.
\citep{exec19_javaid_reference}

Even marginal intervention risks may not settle an individual counterfactual.
In a second exact construction, let $A$ and $U$ be independent fair binary
variables. Compare potential outcomes $(Y(0),Y(1))=(U,U)$ with
$(Y(0),Y(1))=(U,1-U)$. Both generate the uniform observed distribution over
$(A,Y)$ and both have intervention risks $1/2$ for either approval value. For
an observed case with $A=0,Y=0$, however,
\begin{equation}
 P\{Y(1)=1\mid A=0,Y=0\}=
 \begin{cases}0,&(Y(0),Y(1))=(U,U),\\
 1,&(Y(0),Y(1))=(U,1-U).
 \end{cases}
 \label{exec19:javaid-coupling}
\end{equation}
The missing information concerns the joint coupling of potential outcomes,
which the two marginal intervention distributions do not supply. An agent
translating the business question must specify whether it asks for a population
intervention, conditional intervention or cross-world event. Substituting a
Bayesian conditional for the last answers another question.
\citep{exec19_javaid_reference}

\subsection{The selected updater has a narrower evidence contract}
\label{exec19:javaid-evidence}

The selected \texttt{BayesUpdater} documentation describes a joint marginal
conditional on evidence. The \texttt{Identifiability} implementation instead
comments its target as $P_t(s)$, with intervention values obtained from its
evidence object. \texttt{Evidence.getVariablesInEvidence} collects variables
with a single allowed category without checking manipulation flags.
\texttt{Identifiability.getJointMarginal} uses those variables as interventions
and explicitly assumes they are manipulated. Its separate graph-construction
method removes incoming edges only where manipulation flags are true.
Ordinary and manipulated evidence are therefore not interchangeable merely
because they share a container. A first proposed profile should admit validated
singleton interventions and reject mixed evidence until that mismatch is
resolved by a reviewed interface and native tests.~\citep{exec19_javaid_source}

There is also a subtype obligation. The public constructor accepts
\texttt{BayesIm}, but the no-evidence path and Q-expression evaluator cast it
to \texttt{MlBayesImObs} to access the observed joint-probability table. A
generic Bayesian-model object is insufficient evidence that the call will work.
The wrapper must preserve the observed-variable/category mapping and validate
the representation, finite probabilities, normalization and support. The
selected probability setter permits NaN; the Q-expression evaluator computes
a reciprocal denominator without a local zero guard. These observations
require numerical and support checks. They do not prove which end-to-end
native inputs produce a nonfinite result.~\citep{exec19_javaid_source}

Failures need distinct types. A recursive identification failure returns
\texttt{null}; the public joint-marginal path maps that branch to $-1$.
Internal consistency failures throw exceptions, whereas contradictory target
and intervention assignments return the valid probability zero and removing
all target variables returns one. The updated-model getter itself returns
\texttt{null}. The inspected GUI wrapper formats numeric marginals but does
not supply a typed identification certificate in that path. Clamping $-1$ to
zero would fabricate a probability; treating a null model as a zero effect
would confuse object availability with causal meaning. Signed treatment
contrasts are a separate type and can legitimately be negative, so a universal
negative-means-failure rule would also be wrong.~\citep{exec19_javaid_source}

\subsection{An effect list is not a confidence set}
\label{exec19:javaid-candidates}

The IDA-related interfaces need separate graph and sampling contracts.
The \texttt{PdagPagIda} documentation mentions PAGs, but its constructor checks
for a legal DAG or PDAG and continuous data. A nearby comment says that the
rPAG-IDA route is disabled because it still needs work. A class name, dormant
branch or broader comment cannot admit a PAG-only input. Fully oriented graphs
may satisfy more than one graph predicate; the observation is the actual guard,
not a categorical claim about every object labeled a PAG.
\citep{exec19_javaid_source}

Candidate failures complicate interpretation. In a selected optimal-adjustment
branch, an exception can append zero and continue. An outer handler instead
logs an error and omits a candidate. A returned zero or shortened list cannot
automatically certify a structural zero or exhaustive graph-compatible effect
enumeration. A proposed adapter should retain every attempted candidate with
its graph, adjustment set, regression status and exclusion reason, and refuse
scientific publication when failures collapse into values. This source-level
observation is neither a native reproduction nor a claim that every call fails.
\citep{exec19_javaid_source}

The distance functions show another distinction. For candidates $\{-2,2\}$
and target zero, finite-set distance is two, whereas distance to the interval
hull $[-2,2]$ is zero. The selected code uses the extreme range for multiple
candidates. For an empty list it returns the target value itself, so target
minus one yields minus one rather than a valid nonnegative distance. These
source-branch examples require an explicit empty status and an accurately named
metric. Neither an effect list nor its hull is a calibrated confidence region
without an argument for graph uncertainty, sampling error and observation law.
\citep{exec19_javaid_source,exec19_javaid_reference}

\subsection{A proposal route with explicit unresolved obligations}

The proposed Java schema separates graph and probability inputs, singleton
interventions, ordinary evidence, query variables and cross-world conditioning.
Its output distinguishes a completed identified query, a reviewed
nonidentification certificate, unknown scope, unsupported input, and execution
or numerical failure. The source's $-1$ branch must not be upgraded to a
universal nonidentification certificate before the algorithm's assumptions and
proof are reviewed. The original-language engine remains a candidate; no port,
native qualification or distribution permission follows from this comparison.

The remaining public research must inspect the named algorithm original,
graph-legality and probability-storage helpers, and other Java candidates under
a recorded search boundary. It must establish whether conditional ID, general
counterfactual or partial-identification queries have actual supported
interfaces, rather than infer absence from selected files. A later admitted
native prototype can test subtypes, manipulation flags, failure sentinels, zero
support, nonfinite values, skipped candidates and graph guards against the
exact references. R and Python symbolic engines are distinct comparators with
their own graph and failure contracts. These tasks can use public sources and
synthetic laws without requiring authentic bank data to clarify the proposal.


% SOURCE docs/release/2026-10-29/source/docs/writing-repair/2026-10-29/src/434c2f8ba0a0a8d2.tex CHAR 0
\section{Julia and JavaScript: preserving the meaning of an identification request}
\label{exec19:otherid-scope}

A bank deciding whether to make a lending-support offer must distinguish a
graphical search from an answer about repayment. A set of variables may justify
one adjustment formula without providing its numerical value. Conversely,
failure to find such a set need not rule out a different identifying formula.
An interface that turns both outcomes into an empty array loses information
precisely where a business user needs an explanation. Julia's
\texttt{CausalInference.jl} and the JavaScript core of DAGitty expose useful
graph operations, but their selected interfaces require different conversions.
The inspected Julia commit of 10 August 2026 declares version 0.19.4; the
JavaScript revision is dated 6 December 2024. These immutable source identities
define the comparison, rather than an assumed correspondence with an installed
release.~\citep{exec19_otherid_julia,exec19_otherid_dagitty}

\subsection{The graph operation, the observable variables and the target}
\label{exec19:otherid-graph}
Write $A$ for the binary offer, $M$ for a subsequent customer response and $Y$
for a binary repayment outcome. An unobserved factor $U$ may influence both
offer assignment and repayment. A graphical intervention deletes the incoming
arrows to $A$; evaluating $P(Y\mid\operatorname{do}(A=a))$ additionally needs
a model or an identified functional of an observed law. The selected Julia
\texttt{do!} method performs the former operation: it removes incoming edges
and returns the modified graph. Its name alone therefore cannot justify a
numerical intervention answer.~\citep{exec19_otherid_julia}

The inspected adjustment searches expose a mandatory set $I$ and an allowed
set $R$. Their intended result is a set $Z$ satisfying $I\subseteq Z\subseteq R$
and the selected graphical criterion. The default $R$ contains all nonendpoint
vertices. This is convenient for a fully observed graph, but a banking graph
often contains expert-specified unobserved variables. Such vertices must be
excluded from the permitted controls: a mathematically useful hidden variable
is not an available data column. The package's example explicitly removes two
unobserved vertices from $R$, and its selected tests include the corresponding
restricted enumeration. These are inspected examples and test assertions, not
tests executed in this proposal.~\citep{exec19_otherid_julia}

For the proposed boundary, observed status is consequently part of the request,
not an optional display annotation. A request must also retain the outcome
horizon, whether $M$ is pre-action or post-action, and the distinction between a
total, direct or path-specific effect. A proposed control set that contains an
unavailable variable is refused before estimation. An analyst cannot silently
replace a missing pre-action cause by a post-action response and retain the
original estimand. The following exact construction makes the monetary cost of
such a substitution explicit.~\citep{exec19_otherid_reference}

\subsection{A positive-law lending example with three different answers}
\label{exec19:otherid-law}
Consider the project-authored structural law
\begin{align}
 U&\sim\operatorname{Bernoulli}(1/2), &
 P(A=1\mid U)&=\tfrac14+\tfrac12U,\nonumber\\
 P(M=1\mid A)&=\tfrac14+\tfrac12A, &
 P(Y=1\mid M,U)&=\tfrac18+\tfrac12M+\tfrac14U.
 \label{exec19:otherid-structural}
\end{align}
The structural randomizations are independent. The graph is
$U\to A$, $U\to Y$, $A\to M\to Y$; it has no direct $A\to Y$ arrow
and no hidden common cause of $A$ and $M$. Every joint cell has positive
probability. This is a stipulated mechanism, not a fitted loan portfolio.
It supplies a transparent exact reference against which a later adapter can
be compared.~\citep{exec19_otherid_reference}

Under intervention, $U$ retains its fair marginal distribution while the law
of $M$ is set by $a$. Thus direct structural enumeration gives
\begin{equation}
 r(a)=P(Y=1\mid\operatorname{do}(A=a))
 =\tfrac18+\tfrac12(\tfrac14+\tfrac12a)+\tfrac14\tfrac12
 =\tfrac38+\tfrac14a.
 \label{exec19:otherid-risks}
\end{equation}
The interventional risks are $3/8$ and $5/8$, with effect $1/4$. In the
observed law, Bayes' rule instead gives $P(U=1\mid A=a)=1/4+a/2$. Hence
$P(Y=1\mid A=a)=5/16+3a/8$: the observed risks are $5/16$ and $11/16$,
with contrast $3/8$. The observed comparison mixes the offer response with
selection on $U$.~\citep{exec19_otherid_reference}

The frontdoor expression evaluated from this law is
\begin{equation}
 f(a)=\sum_m P(m\mid a)\sum_b P(Y=1\mid b,m)P(b).
 \label{exec19:otherid-frontdoor}
\end{equation}
Here $U$ and $M$ are conditionally independent given $A$. Therefore
$P(Y=1\mid b,m)=3/16+m/2+b/8$ and $P(b)=1/2$. The inner average is
$1/4+m/2$, so $f(a)=3/8+a/4=r(a)$. This verifies the expression for
the stipulated law. It is not a proof of a general frontdoor theorem, a
completeness claim for the selected software, or an upstream execution. The
proof originals named by the source remain separate inspection obligations.
\citep{exec19_otherid_reference,exec19_otherid_julia}

Using $M$ as an ordinary standardization control gives a third answer. Its
marginal success probability is $P(M=1)=1/2$, so
\begin{equation}
 h(a)=\sum_m P(Y=1\mid a,m)P(m)=\tfrac7{16}+\tfrac18a.
 \label{exec19:otherid-standardization}
\end{equation}
The resulting contrast is $1/8$, not $1/4$. It is neither the total-effect
answer nor a justified direct-effect answer merely because the mediator has
been included as a control. In the graph, the empty set leaves
$A\leftarrow U\to Y$ open, and the only observed nonendpoint candidate $M$
does not block that backdoor path. The absence of a permitted backdoor set
therefore coexists, in this example, with the successful expression
\eqref{exec19:otherid-frontdoor}.~\citep{exec19_otherid_reference}

For an assumed incremental success payoff of USD 100 and an offer cost of
USD 20, the three contrasts imply
\begin{equation}
 100(1/4)-20=5,\qquad
 100(3/8)-20=35/2,\qquad
 100(1/8)-20=-15/2.
 \label{exec19:otherid-money}
\end{equation}
The same positive observed law can thus tempt a user to overstate the proposed
offer value as USD 17.50 or to reject it at USD $-7.50$, whereas the stipulated
intervention has value USD 5. These are assumed-money consequences of query
semantics, not observed profits, estimated confidence intervals or an investment
recommendation. Sampling uncertainty and the credibility of the graph would
still require separate analysis.~\citep{exec19_otherid_reference}

\subsection{A failed search needs a reason and a completion status}
\label{exec19:otherid-status}
The Julia find methods return a set or \texttt{false}; their enumeration
iterator returns \texttt{nothing} when exhausted and declares an unknown
iterator size. One empty set is a valid candidate distinct from an enumeration
that emits no sets. The inspected \texttt{find\_backdoor\_adjustment} path
passes the result of \texttt{find\_covariate\_adjustment} directly into
\texttt{setdiff}, without an explicit local guard for \texttt{false}. This
is a specific native conformance question, not a reproduced exception. A
future wrapper must preserve an error as an error rather than converting it
to scientific nonidentification.~\citep{exec19_otherid_julia}

DAGitty illustrates why a single Boolean or list cannot carry the full
contract. Its selected total-adjustment path admits DAG, PDAG, MAG and PAG
tags, whereas its direct-effect check restricts the graph to a DAG and rejects
selection nodes. Missing source or target markings, an invalid adjustment set,
and selected failures of adjustment or selection conditions can produce
\texttt{false} or an empty list through different interfaces. Its total-effect
enumeration retains latent and selection annotations and accepts a maximum
result count. The selected MAG/PAG validation branches themselves retain
comments that proper validation is unfinished. Acceptance of a graph tag must
therefore not be promoted into validation of every scientific premise.
\citep{exec19_otherid_dagitty}

The proposed shared result has separate fields for request validity, graph
profile, search criterion, emitted solutions, completion and scientific
conclusion. A budget-limited result preserves every emitted set and reports
incomplete or unknown completion. Reaching a cap does not by itself prove that
another solution exists; nor does it certify exhaustiveness. A valid completed
adjustment search can support ``no route under this adjustment criterion.''
A general ``not identified'' conclusion instead needs an applicable reviewed
certificate. A missing endpoint, unsupported graph, implementation exception
or unverified proof premise remains a different outcome. None defaults to a
zero effect.~\citep{exec19_otherid_julia,exec19_otherid_dagitty,exec19_otherid_reference}

\subsection{A narrow adoption route and the questions it leaves open}
\label{exec19:otherid-route}
A first prototype can use a narrow declared DAG profile with explicit observed
vertices, disjoint query sets, validated $I\subseteq R$, and an exact
enumeration oracle on small graphs. It should test a valid empty adjustment
set, no permitted adjustment set with a separate frontdoor route, an invalid
mandatory set, a hidden candidate, missing endpoints, invalid graph marks and
a deliberately capped enumeration. Compare sets independently of order and
retain completion separately from the count. Execute pinned upstream tests only
after recording the native environment and distribution; their presence in a
repository is useful evidence of intended behavior, not a passing local run.
\citep{exec19_otherid_julia,exec19_otherid_dagitty}

The Julia metadata supplies dependency ranges and an MIT notice, rather than
a fully resolved distribution. The JavaScript candidate retains the component
and dependency review obligations already identified in its reuse comparison.
An adapter boundary is a candidate for evaluation, while a wholesale port
would additionally require maintained correspondence with the original
algorithms and tests. A source-level convenience is not sufficient reason to
incur that maintenance burden. No executable integration permission is awarded
by this comparison.~\citep{exec19_otherid_julia,exec19_otherid_dagitty}

Three unresolved tasks remain distinct. Broader language and package discovery
is needed before claiming coverage. The source-linked adjustment, frontdoor
and path-specific originals require genuine PDF acquisition and inspection
before their general guarantees are cited. Finally, native failure behavior,
graph validation and distribution correspondence need a separately reviewed
prototype. The recanting-witness helper concerns a specified path-specific
question under complete observations; its name is not evidence of arbitrary
latent or cross-world identification. Public graph fixtures and stipulated
laws can address these prototype questions before any expensive bank trial.
Authentic assignment knowledge, outcome truth and institutional acceptance
remain later evidence.~\citep{exec19_otherid_julia,exec19_otherid_dagitty,exec19_otherid_reference}


% SOURCE docs/release/2026-10-29/source/docs/writing-repair/2026-10-29/src/f7dc2abb67a65b90.tex CHAR 0
\section{Counterfactual comparisons in C++ libraries}
\label{exec19:cxxid-scope}
A lending analyst may ask whether an approval changes repayment, while a
servicing analyst asks whether a borrower who failed under the existing policy
would have succeeded under another action. These questions require different
information. A native implementation can make a symbolic calculation faster,
but its language does not turn an observational conditional distribution into
an intervention law or select a coupling between a person's potential outcomes.
The proposed product therefore keeps graph identification, numerical evaluation
and counterfactual model specification as separate contracts.

The selected aGrUM 3.2.0 revision is dated 6 October 2026. Twelve authenticated
implementation, wrapper and test originals extend the earlier dispatch review.
The build declarations require C++20 and CMake 3.18; the selected root license
statement is LGPL-3-or-later or MIT, subject to a separate component review.
These are historical source declarations, not a report that a native build,
Python installation or bank interface has passed. The R-backed C++ search core
of dosearch supplies a substantive alternative, while other native candidates
and partial-identification engines remain unresolved. The bounded comparison
cannot establish complete C++ market coverage.
\citep{exec19_cxx_identification_source,exec18_rid_dosearch}

\subsection{Inspecting recursion without mistaking it for a correctness proof}
\label{exec19:cxxid-recursion}
The native do-calculus source constructs backdoor and frontdoor expressions,
then provides a general recursive identification path. Its branches restrict
the graph to outcome ancestors, add irrelevant remaining variables to the
intervention, partition observed nodes connected through latent co-children,
and multiply or marginalize the resulting kernels. A single remaining full
confounded component triggers a hedge exception. The induced submodel keeps
observed identifiers and retains a latent descriptor only when it still has at
least two retained children. These observations identify the algorithmic
objects that a comparison must preserve; they are not a proof that every legal
input reaches the mathematically intended branch.
\citep{exec19_cxx_identification_source}

Two details deserve explicit qualification before reuse. The conditional-query
path builds one modified graph after removing incoming edges to intervention
nodes and outgoing edges from every conditioning node, then tests candidate
promotions within that graph. A comparison must establish the appropriate
multi-variable graph surgery and the conditional query after each promotion.
Separately, when an accumulated kernel is present, the first topological
factor is constructed as a joint-probability AST rather than the numerator
expression just calculated from that kernel. Whether that branch is reachable
with a distinct marginal, or is equivalent under the implementation's invariant,
requires a proof-aligned fixture. Neither observation alone warrants announcing
a software defect. Both are precise obligations that a root README or successful
simple adjustment example would miss.
\citep{exec19_cxx_identification_source}

The numerical layer has its own boundary. A formula's identified predicate
checks whether an expression root exists. Evaluation then requires an
observational Bayesian network, so a structure-only answer and an evaluable
answer are different states. Division of expression trees delegates to tensor
division. The earlier selected files do not settle zero-denominator support
handling, finite arithmetic or an uncertainty statement. The more detailed
comparison in Section~\ref{sec:exec24-agrum-reuse} resolves a bounded part of
that source question: selected tensor-division branches introduce no support-zero
refusal, ordinary zero-sum normalization leaves the tensor unchanged, and CPT
normalization refuses a zero-sum conditioning row. Those static findings specify
native fixtures; they do not establish universal runtime behavior or statistical
uncertainty.~\citep{exec24_agrum_source,exec24_agrum_choice}
The Python convenience wrapper
retains the model on the returned tensor and maps an unidentified formula to a
\texttt{None} numerical answer; complete exception and lifetime conversions
still need direct inspection and execution. Preserve an expression with its
graph, variable domains and support conditions instead of returning an
unqualified floating-point success.
\citep{exec19_cxx_identification_source}

For example, let approval $A$ and repayment $Y$ be binary. In one authored
lending world, $U\sim\operatorname{Bernoulli}(1/2)$, $A=U$ and $Y=U$;
in another, $A\sim\operatorname{Bernoulli}(1/2)$ and $Y=A$. Both observed laws
put mass $1/2$ on each of $(A,Y)=(0,0)$ and $(1,1)$. Replacing the approval
equation gives an average intervention contrast of zero in the first world
and one in the second. Deterministic assignment in the first world also means
that the observed law does not provide all conditional treatment comparisons.
The example is a supplied structural construction, not a discovered bank graph.
Its purpose is to prevent an association or unsupported engine return from
being presented as identified repayment benefit.
\citep{exec19_cxx_identification_reference}

\subsection{Why individual root updates need a joint abduction contract}
\label{exec19:cxxid-abduction}
The selected counterfactual builder starts from the observational Bayesian
network, collects parentless nodes excluding interventions and declared latents,
conditions the original network on the observed profile, and fills each selected
root's conditional-probability table with its individual posterior. It then
calculates an intervention in the reconstructed model. The constructor used
for reconstruction does not copy the separate original latent descriptors or
spurious-arc declarations. Accordingly, graph identity also belongs in the
conformance test. More fundamentally, replacing several root marginals does
not in general retain their posterior joint dependence.
\citep{exec19_cxx_identification_source}

Consider a completely declared servicing example. Before an action, roots
$U$ and $V$ are independent fair bits; $A$ is an independent fair action.
Let success under action $a$ have one of two structural response functions:
\begin{align}
 Y_{\mathrm{stable}}(a)&=\mathbf 1\{U=V\},\nonumber\\
 Y_{\mathrm{switch}}(0)&=\mathbf 1\{U=V\},\qquad
 Y_{\mathrm{switch}}(1)=\mathbf 1\{U\ne V\}.
 \label{exec19:cxxid-coupling}
\end{align}
Every observed pair $(A,Y)$ has probability $1/4$ in both worlds; each world
also has $P\{Y(a)=1\}=1/2$ for both actions. Thus even the two marginal
intervention laws do not choose between these cross-world response functions.
For the profile $A=0,Y=1$, however, the roots must agree. The joint posterior
puts mass $1/2$ on $(U,V)=(0,0)$ and $1/2$ on $(1,1)$, while each separate
root posterior remains a fair bit. Consequently,
\begin{align}
 P\{Y_{\mathrm{stable}}(1)=1\mid A=0,Y=1\}&=1,\nonumber\\
 P\{Y_{\mathrm{switch}}(1)=1\mid A=0,Y=1\}&=0,\nonumber\\
 \sum_{u,v}\mathbf 1\{u=v\}P(u\mid A=0,Y=1)P(v\mid A=0,Y=1)&=\tfrac12.
 \label{exec19:cxxid-joint}
\end{align}
The last expression is the result of imposing separate posterior root
marginals; it is not the first world's joint conditional probability. In the
switch world the same product reconstruction also gives $1/2$ for inequality.
The exact finite enumeration tests this semantic distinction without executing
an upstream package. A native run is still required to establish the behavior
of the installed representation, graph and wrapper on this fixture.
\citep{exec19_cxx_identification_reference}

If the declared payoff is USD 100 for success less a USD 20 action cost, the
three respective success probabilities imply USD 80, USD $-20$ and USD 30.
These are conditional values for the specified profile and action, not an
incremental return against a measured current workflow. They are assumed
monetary consequences of a semantic choice. A product cannot substitute the
middle probability $1/2$ because it looks numerically stable: it must preserve
the joint abduction distribution, state a sufficient factorization condition,
return a justified set over admissible couplings, or refuse the requested
individual answer. Additional marginal intervention observations alone do not
resolve this constructed coupling ambiguity.
\citep{exec19_cxx_identification_reference}

\subsection{A reusable boundary and a concrete qualification route}
\label{exec19:cxxid-contract}
The inspected native tests are useful starting specifications. They include
manual backdoor/frontdoor expression comparisons, selected direct-do-calculus
numerical comparisons, a hedge evaluation exception, ID/IDC smoke tests and a
salary counterfactual fixture. The selected conditional smoke tests verify that
an expression builds and evaluates; they do not independently establish its
multi-variable causal meaning. A parity check against the same builder also
cannot replace a separate structural oracle. None of these upstream tests was
run in this source review, and the coupled-root example is a new authored
reference, not a claim about existing upstream coverage.
\citep{exec19_cxx_identification_source,exec19_cxx_identification_reference}

For dosearch, the already pinned R/C++ source makes a different tradeoff. Its
native entry accepts directed and bidirected graph information, a collection
of known distributions, a target and search options through Rcpp. The loop can
return when its time budget expires; the returned identifiable flag depends
on whether a target distribution was found. The inspected result does not
separately certify search exhaustion. Therefore preserve a discovered formula
and derivation when available, but retain unknown or time-limited status when
search failure has not been proved to imply nonidentification in the stated
problem class. Start reuse through the supported R boundary; a stand-alone C++
port would additionally require extraction of the Rcpp interface, environment,
representation and rights obligations. A shared implementation language is
insufficient evidence that the two engines deliver an interchangeable service.
\citep{exec18_rid_dosearch}

The proposed conformance program first fixes the query, graph class, observed
or experimental input law, domains, support and information clock. It then
compares direct and shortcut identification on independently enumerated small
models, tests multi-variable conditioning and nested kernels, exercises
nonidentification separately from unsupported methods, and tests numeric
availability and support-zero behavior. Counterfactual tests additionally
retain structural response functions and joint posterior dependence; a graph
with ordinary conditional tables is not silently promoted to a unique
cross-world model. Finally, the same signed query and typed result cross the
native, Python and process boundaries, with ownership and exceptions checked.
Retain aGrUM and dosearch as conditional comparators and defer executable
admission until those trials pass in their declared regimes. Other candidates,
original proof alignment, full component rights and general partial
identification remain open; this bounded source review does not close them.
\citep{exec19_cxx_identification_source,exec18_rid_dosearch,exec19_cxx_identification_reference}


% SOURCE docs/release/2026-10-29/source/docs/writing-repair/2026-10-29/src/ceb362b360a2a9ec.tex CHAR 0
\section{Ananke reuse: graph identification, model oracles and effect scales}
\label{sec:exec24-ananke-reuse}

A bank translating a collections decision into a causal query needs to know
whether a program identifies an intervention distribution, evaluates an
already supplied structural model, or estimates a particular contrast.
Ananke offers distinct interfaces for these tasks. Their proximity within
one distribution is useful for research, but does not make their inputs,
guarantees or results interchangeable. The proposed reuse decision is to
retain the selected source as a reference and consider an isolated
comparator after explicit scientific and engineering qualification.
\citep{exec24_ananke_source,exec24_ananke_choice}

The inspected distribution is \texttt{ananke-causal} 0.5.0, uploaded on
5 December 2023. Its preserved wheel and selected extracted members are
authenticated byte for byte. This is a bounded historical comparison,
not a claim about the latest release or current support. The metadata
declares Apache 2.0 and also an Other/Proprietary classifier; two matching
Apache-text files are retained. The conflicting declaration, component
notices and dependency rights remain distribution questions to resolve.
Authenticity of bytes does not settle those questions.
\citep{exec24_ananke_source}

\subsection{Identification requires a declared graph and information set}
For named treatment variables $A$ and outcomes $Y$, the selected
\texttt{OneLineID} implementation fixes the treatment vertices, constructs
the relevant outcome ancestors and checks district fixability. Its
Boolean result concerns the supplied graph and query. A causal graph
elicited by an analyst or an LLM remains an assumption-bearing scientific
input; the algorithm does not validate that graph against banking truth.
\citep{exec24_ananke_source,exec24_ananke_choice}

The functional interface constructs a display string and prints fixing
orders. Some variable lists are concatenated without separators. As a
simple authored encoding example, the token arrays $(A,BC)$ and $(AB,C)$
both concatenate to ABC. This does not demonstrate that the algorithm
returns the same functional for different valid queries. It shows why a
display representation cannot serve as a lossless interchange contract.
The proposed CLI/MCP boundary needs named-variable arrays and a typed
expression tree preserving sums, products, kernels, fixing order and
available distributions. It must not evaluate returned display text.
\citep{exec24_ananke_source,exec24_ananke_choice}

The experimental extensions have additional input restrictions. The
selected ancestral-experiment and generalized-experiment interfaces can
raise \texttt{NotIdentifiedError} when their experiment preconditions are
not satisfied. The same exception class therefore cannot be translated
mechanically into proof of nonidentification. A worker must distinguish
invalid input, unsupported information, scientific nonidentification,
resource termination and engine failure, retaining the original condition.
Independent graph preflight also remains required; a source comment
requesting correctness review is not itself an observed package failure.
\citep{exec24_ananke_source}

\subsection{A true-model oracle and missing-data query answer different questions}
The discrete-model oracle explicitly performs interventions in a supplied
true model without considering identification. Conditional results divide
an intervened joint distribution by an intervened conditioning probability.
This is useful for public fixtures whose data-generating model is known.
It cannot turn observational ambiguity into knowledge of the correct model,
and normalization requires a nonzero conditioning probability. A banking
result must state which information was supplied and which was inferred.
\citep{exec24_ananke_source,exec24_ananke_choice}

The missing-data full-law routine uses the naming conventions
\texttt{R\_} and \texttt{X\_} and constructs an associated variable name
from a split indicator name. In the selected source, an indicator called
\texttt{R\_credit\_score} leads that string operation to
\texttt{X\_credit}, rather than \texttt{X\_credit\_score}.
This authored naming witness does not establish a native erroneous result.
It establishes a conversion obligation: expose a checked bijection between
business field identifiers, missingness indicators and internal graph
names, or refuse an unsupported name. A full-law identification check
also does not itself provide arbitrary individual counterfactual inference.
\citep{exec24_ananke_source,exec24_ananke_choice}

\subsection{The requested effect scale must match every returned interval}
The average-effect interface selects estimators according to a-, p- and
nested-fixability conditions. Its reporting method recognizes variables
whose observed values are subsets of zero and one as binary, adds a column
named \texttt{ones}, and obtains fitted means at assignments zero and one.
These choices require an explicit schema for treatment values, missingness,
reserved names and outcome type. They are not a general multi-valued
treatment contract.~\citep{exec24_ananke_source}

For binary outcomes the default point calculation returns a log odds
ratio. When \texttt{report\_log\_odds=False}, the point calculation instead
uses the difference of fitted means. The inspected bootstrap branch tests
binary outcome type without also testing that reporting flag; it then
accumulates log odds ratios before taking quantiles. Thus a specific
binary-outcome call with the flag false has different point and replicate
reporting branches. This is a source-traced conditional inconsistency,
not a claim that a native trial was performed or that all uses are affected.
\citep{exec24_ananke_source}

An authored finite witness isolates the units. Suppose fitted means are
$p_0=1/5$ and $p_1=2/5$. The risk difference is
\[
 p_1-p_0=1/5,
\]
whereas the odds ratio and its logarithm are
\[
 \frac{p_1/(1-p_1)}{p_0/(1-p_0)}=8/3,
 \qquad \log(8/3)\simeq0.980829.
\]
If hypothetical replicate means all equal these values, quantiles of
the log odds ratios equal $\log(8/3)$, while the requested difference
point remains $1/5$. Such endpoints cannot be labeled as a risk-difference
interval. The example does not assert that a particular fitted estimator
produces those replicates. It supplies a small reporting oracle for a
future corrected adapter or a deliberately restricted wrapper.
\citep{exec24_ananke_choice}

The same selected method samples rows with replacement and calls ordinary
quantiles. The local loop does not expose a cluster/calendar resampling
contract, a dedicated random-state argument or a per-replicate status
ledger. Finite fitted probabilities at zero or one also need explicit
handling before an odds transformation. A bank panel or dependent servicing
trajectory must not inherit an independent-row uncertainty argument merely
because bootstrap endpoints are available. Native trials must count every
requested replicate, invalid fit, retry and refusal.
\citep{exec24_ananke_source,exec24_ananke_choice}

\subsection{Reuse choices, alternatives and qualification}
The selected metadata constrains NumPy below version 2, pandas below 2 and
statsmodels below 0.14, alongside older pgmpy bounds and JAX components.
Those declarations differ from other inspected candidate environments.
The symbolic identification module also imports SymPy without a direct
declaration in this selected metadata. Transitive provision, dependency
compatibility and a reproducible installed environment remain unverified.
An isolated worker is a proposed means of containing version conflicts,
not evidence that a deployment works.~\citep{exec24_ananke_source}

The proposal retains Ananke as a source reference now. Conditional integration
as a narrow graph-identification comparator may be useful after a frozen
query contract, environment and rights review. Direct adoption of an
unrestricted identification-and-estimation service is deferred. A full
source port is also deferred because it transfers graph correctness,
symbolic semantics and numerical maintenance without a measured benefit.
DoWhy's explicit model/identify/estimate separation, y0's typed symbolic
representation and restricted pgmpy adjustment remain inherited alternatives;
their own source limits are not removed by this comparison.
\citep{exec24_ananke_choice}

The scientific owner must fix the graph, available laws, target and
sampling assumptions. The adapter owner must preserve node identities,
expression structure, effect scales and error categories. The platform
and rights owner must resolve the distribution and its complete dependency
environment. The evaluation owner must freeze public graph/query fixtures,
numerical and refusal oracles, and acceptance criteria before testing.
An illustrative 3,000 USD workflow over 150 requests, 125 finite results
and 100 accepted results costs 20, 24 and 30 USD respectively. Those are
authored prices; measured engineering, review, all-request execution,
storage and maintenance costs remain unknown.
\citep{exec24_ananke_choice}

Concrete tests should cover valid and invalid graphs, empty valid adjustment
sets, nonidentified and unsupported queries, experimental preconditions,
name round trips, zero normalization denominators, binary reporting flags,
probability boundaries and dependent resampling. Matched known-law public
fixtures can establish whether a prototype merits costly bank experiments.
They cannot supply unknown assignment truth or identify every cross-world
quantity without added information. Native, independent and final
acceptance remain open.~\citep{exec24_ananke_choice}


% SOURCE docs/release/2026-10-29/source/docs/writing-repair/2026-10-29/src/21a0fcecc8d4f6ff.tex CHAR 0
\section{DoWhy GCM reuse: supplied mechanisms, counterfactual information and uncertainty}
\label{sec:exec24-gcm-reuse}

A bank asking how a servicing intervention changes default risk may also
ask whether a particular customer would have avoided default. These are
different information requests. DoWhy's graphical causal model (GCM)
interfaces offer useful ways to represent mechanisms and propagate
interventions, but they require a supplied causal interpretation. They do
not inherit approval from a separate review of DoWhy's identification API.
The proposed decision retains selected GCM sources as a reference for a
narrow future worker.~\citep{exec24_gcm_source,exec24_gcm_choice}

Twelve selected original files authenticate to the repository revision of
4 October 2026: ten GCM and graph files and the retained project and MIT
license declarations. The project uses dynamic versioning; its 0.0.0
placeholder is not a release version. The declarations require Python
at least 3.9 and below 3.14, NumPy above 2.0 and version-dependent SciPy
and NetworkX constraints. A complete distribution, numerical environment,
transitive-rights review and tested import path remain outstanding.
No downloaded package or test suite is executed here.
\citep{exec24_gcm_source}

\subsection{A supplied graph and fitted mechanisms are separate scientific commitments}
For a directed acyclic graph, write a proposed structural system as
\[
 X_j=f_j(X_{\operatorname{pa}(j)},N_j).
\]
The parent set, measurement clock, exogenous-noise law and invariance
under the proposed action determine the interpretation. For the selected
independently sampled mechanisms, a scientific model must justify the
corresponding factorization and absence of omitted shared noise. A conditional
predictor for $X_j$ is not alone evidence that its inputs are causes or
that its residual is independent of them. Unrecorded branch policy or
common calendar shocks can invalidate that interpretation even when
prediction is accurate.~\citep{exec24_gcm_choice}

The selected interfaces distinguish stochastic conditional mechanisms,
functional mechanisms and mechanisms invertible with respect to noise.
The additive-noise path fits a predictor and a residual distribution;
the post-nonlinear path first inverts an outer transformation, fits a
predictor, and reconstructs residual noise. These steps compute a model
under declared restrictions. They do not demonstrate those restrictions
for banking data.~\citep{exec24_gcm_source}

Parent columns are selected in sorted node order. The fitted model records
the ordered parents and checks them before causal queries, detecting a
changed parent set. This is useful structural consistency checking; it
does not establish variable identity, units, observation timing or causal
truth. The proposed worker uses stable typed node identifiers, preserves
the fitted ordering, and tests asymmetric parent permutations and changed
graphs. A graph clone deliberately has untrained cloned mechanisms, so
clone and fitted snapshot are different lifecycle states.
\citep{exec24_gcm_source,exec24_gcm_choice}

Fitting filters rows with missing target values for each node. Parent
missingness remains a matter for the chosen predictor. Consequently
different mechanisms can be fitted on different selected rows. The
automatic-assignment interface rejects missing data by default and offers
an explicitly experimental numerical-missingness option. A bank workflow
needs per-mechanism row masks, a missingness argument and tests that its
joint model still describes the intended population. Local deletion is
not an identification argument for missing outcomes.
\citep{exec24_gcm_source,exec24_gcm_choice}

Automatic assignment chooses an empirical distribution for a root and
selected regression or classification mechanisms for other nodes. Its
selection uses shuffled predictive folds and MSE or log loss. It assigns
mechanisms to the supplied graph; it does not discover the graph or prove
noise independence. The selected scorer suppresses convergence and runtime
warnings and can skip classification folds with one class. A proposed
evaluation retains effective fold counts and diagnostics, and refuses
selection when every candidate has a nonfinite score. Calendar and
customer splits need an explicit alternative to shuffled row folds.
\citep{exec24_gcm_source,exec24_gcm_choice}

An LLM may propose variables, edges, transformations and mechanism classes
from a business description. Those proposals should carry sources and
alternative diagrams, then receive human approval before fitting. Predictive
model selection and an agent's confidence cannot replace the intervention,
clock or latent-variable assumptions.~\citep{exec24_gcm_choice}

\subsection{Intervention sampling does not retain an individual's noise}
The intervention interface accepts either observed baseline rows or a
requested generated sample count, not both. It copies the baseline data,
updates affected nodes in topological order and draws new conditional
values for affected non-root nodes before applying the action. An atomic
action replaces a value; a value-dependent transformation acts on the
newly generated pre-intervention value at that node. Unaffected baseline
columns can remain observed. This is a specified population or baseline
sampling procedure, not individual abduction.~\citep{exec24_gcm_source}

The action evaluator applies the provided function to each scalar and
checks the resulting shape. A proposed CLI or MCP interface should accept
a reviewed declarative action grammar, such as fixed-value replacement
or a bounded affine transform, rather than transmitting arbitrary code.
The contract also records domains, units, interventions on descendants
and the model interpretation of each transformation.
\citep{exec24_gcm_source,exec24_gcm_choice}

The counterfactual path instead propagates supplied noise or reconstructs
noise from observed rows using invertible mechanisms. For an authored
invariant system $Y=2A+N$, an observed account with $A=1,Y=5$ has
reconstructed $N=3$. Holding that noise fixed gives $Y(0)=3$.
A population intervention under the additional assumption $E[N]=0$
instead has $E[Y(0)]=0$. The two answers are consistent because the
targets differ. Neither calculation validates the supplied equation
against the bank's actual decision process.~\citep{exec24_gcm_choice}

Categorical functional mechanisms use a uniform noise coordinate and
cumulative class probabilities. An observed category generally determines
an interval of that noise, not a unique value. The selected abduction
helper explicitly rejects a noninvertible mechanism; passing noise
directly is a different information contract. A worker must retain whether
noise was measured, stipulated or inferred and what assumptions justify
that inference. Inventing a noise value to bypass the refusal would
manufacture a different counterfactual question.
\citep{exec24_gcm_source,exec24_gcm_choice}

An authored binary example makes the information gap precise. Let
$U$ be a fair binary variable and randomize action independently. Model
I has $Y(0)=U,Y(1)=U$; Model II has $Y(0)=U,Y(1)=1-U$.
Both have Bernoulli$(1/2)$ outcome distributions under either action,
the same randomized observed law and average effect zero. Yet among
accounts with $Y(0)=0$, the probability of $Y(1)=1$ is zero in Model I
and one in Model II. Population intervention data alone cannot choose
between these cross-world couplings. The proposed service returns a
model-conditional answer or an identified set when the required additional
structure is absent.~\citep{exec24_gcm_choice}

\subsection{Average-effect event coding and intervention scope must remain explicit}
The selected average-effect helper restricts the model to ancestors of
the requested target, then validates intervention names in that reduced
graph. A request intervening on a disconnected node can therefore be
rejected rather than returned as a model-implied zero. Both are meaningful
states if documented; an adapter must not turn an exception into an
estimated zero.~\citep{exec24_gcm_source,exec24_gcm_choice}

For a binary classifier target, the positive class fixes the event whose
probability is contrasted. The selected API permits an explicit class;
otherwise the second internal class is positive. An authored default-risk
change from $1/5$ to $2/5$ is $+1/5$, while the complementary nondefault
event changes by $-1/5$. For a stipulated 1,000 accounts and 100 USD loss
per default, the former implies 20,000 USD additional expected loss.
Event labels and financial signs must survive the library, CLI and MCP
round trips. These quantities are authored assumptions, not observed
bank effects.~\citep{exec24_gcm_source,exec24_gcm_choice}

The alternative and reference interventions share a baseline table but
invoke conditional sampling separately. A shared table does not imply
shared downstream noise. A reproducibility contract should specify seeds,
mechanism randomness, concurrency and any deliberate common-random-number
construction; Monte Carlo variance must be separated from uncertainty
about the fitted causal model.~\citep{exec24_gcm_source,exec24_gcm_choice}

\subsection{A confidence wrapper inherits the experiment performed by its callable}
The selected generic confidence wrapper repeats a callable with generated
seeds and returns coordinatewise percentile limits and a geometric-median
summary by default. Its default is 20 repetitions. The helper named
bootstrap sampling is an alias for partial application: by itself it
does not resample training data or refit mechanisms. Repeating a query
on a fixed fitted model can measure simulation variation while omitting
sampling, selection, graph and misspecification uncertainty.
\citep{exec24_gcm_source}

The separate fit-and-compute helper clones the model and chooses a
fraction of training rows without replacement, with default fraction
$0.75$. It refits the mechanisms and optionally repeats automatic
assignment; the graph is still the supplied graph. Without-replacement
subsampling and the ordinary size-$n$ bootstrap are different laws.
For authored data $(0,0,2,2)$, all size-three subsets have means
$(2/3,2/3,4/3,4/3)$, whose variance about the empirical mean is $1/9$.
The size-four with-replacement empirical bootstrap mean instead has
variance $1/4$. These are exact conditional finite-population calculations,
not evidence about the coverage of either interval procedure.
\citep{exec24_gcm_source,exec24_gcm_choice}

Subsampling fractions, dependence units, omitted graph uncertainty and
estimator-specific centering or scaling require a justified uncertainty
design. A narrow distribution of repeated simulations cannot settle those
questions. Coordinatewise percentile intervals also do not automatically
supply simultaneous coverage for a portfolio of causal queries.
\citep{exec24_gcm_choice}

Dictionary outputs require a further alignment contract. The selected
wrapper appends each returned value to its key's list, then stacks the
lists for the joint summary. It does not locally require every repetition
to return the same key set. An authored sequence returning key $a$ only
in repetitions $0,2,4$ with values $(0,1,2)$ and key $b$ only in
repetitions $1,3,5$ with values $(0,10,20)$ has equal list lengths but
no joint repetitions. Stacking produces artificial pairs
$(0,0),(1,10),(2,20)$. A proposed worker must preserve repetition IDs,
require a fixed output schema or explicitly represent missing values,
and report requested and valid denominators per query. This is a
conditional source deduction, not an executed package failure.
\citep{exec24_gcm_source,exec24_gcm_choice}

\subsection{Graph falsification is a diagnostic with a defined comparison law}
The selected falsification routine compares local-Markov-condition
violations with those from node-permuted graphs and separately assesses
whether permutations commonly remain in the graph's equivalence class.
Its result distinguishes evaluability, falsifiability and rejection.
A zero-test case is not evaluable; data-subset evaluation is allowed by
default. These states and the exact tested subset should be retained
instead of compressed into a graph-approved Boolean.
\citep{exec24_gcm_source}

This diagnostic does not prove graph truth, rule out all hidden causes
or identify an unrestricted intervention. Its conditional-independence
tests, dependence law, permutation budget and power need their own
qualification. The linked scholarly justification has not been newly
reviewed in this selected-source decision and belongs in its follow-up
register. Source comments are not a substitute for that paper's authentic
PDF and inspected technical argument.~\citep{exec24_gcm_choice}

\subsection{Reference decision, ownership and public qualification}
The immediate choice is reference only. Direct adoption of all GCM paths
is deferred. A narrow isolated integration may follow a frozen model and
action contract, full environment and rights review, and known-law public
fixtures. Porting would transfer responsibility for mechanism fitting,
graph lifecycle, randomness, abduction and uncertainty without a demonstrated
cost advantage. Unsupported causal certainty and invented individual noise
are rejected claims; the selected mechanisms remain useful candidates.
\citep{exec24_gcm_choice}

The scientific owner approves the graph, target, invariance and noise law.
The adapter owner preserves fitted parent identity, action grammar,
positive class, noise provenance, repetition IDs and all refusal states.
Platform and rights owners resolve dependencies, resource limits and
distribution obligations. Evaluation records every requested fit and query,
including missingness, failed selection, noninvertibility, timeouts and
nonfinite results. An authored 5,400 USD total across 270 requests,
225 finite outputs and 180 accepted outputs yields 20, 24 and 30 USD
under the respective denominators. These are illustrative costs.
\citep{exec24_gcm_choice}

Public qualification should compare independent structural oracles with
the pinned worker, including asymmetric graphs, parent reordering,
changed mechanisms, categorical event coding, disconnected interventions,
missing-data masks, individual-versus-population targets and all uncertainty
routes. It should include deliberately indistinguishable cross-world models
to test appropriate refusal. Hidden banking truth and unrestricted dependent
uncertainty have no universal software repair. Execution of the package, independent validation and evaluation in a bank remain outstanding.~\citep{exec24_gcm_choice}


These identification comparisons explain why refusal can be informative.
An unidentified individual response, an unsupported graph class, an exhausted
search and a failed numerical calculation have different remedies. The
worked laws also show when an unresolved individual quantity need not block
an average-value decision, and when a harm constraint makes that same
ambiguity decisive. A shared interface earns its value by keeping those
distinctions available to the decision maker. The next two chapters follow
the result into estimating equations and covariance calculations, where
population weights, nuisance fits and dependence introduce further obligations.

